% concatenated from DDP_Modulus-big.tex
\documentclass[a4paper,11pt]{amsart}
\usepackage[utf8]{inputenc}
\usepackage[english]{babel}
\usepackage{fancyhdr}
%\pagestyle{fancy}
\usepackage{amsmath}
\usepackage{amsfonts,amstext}
\usepackage{amssymb,amsthm,amscd,amsxtra,wasysym,graphicx}
%\usepackage[OT1]{fontenc}
\usepackage{mathrsfs,esint,comment}
\usepackage{tikz-cd}
\usepackage{enumitem,mathtools,dsfont}
 \usepackage{palatino,mathpazo}
 \usepackage{palatino}
\usepackage{charter}
\usepackage{xcolor}
\def\blue{\color{blue}}
\def\orange{\color{orange}}
\def\red{\color{red}}
\def\green{\color{green}}  

\def\Bibtex{{\rm B\kern-.05em{\sc i\kern-.025em b}\kern-0.08em T\kern-.1667em\lower.7ex\hbox{E}\kern-.125emX}}
\hfuzz1pc
 \usepackage[breaklinks=true,bookmarks=false,pagebackref]{hyperref}
 \usepackage{hyperref}
\hypersetup{
     unicode=false,          
    pdftoolbar=true,       
    pdfmenubar=true,       
    pdffitwindow=false,     
    pdfstartview={FitH}, 
    pdftitle={Modulus of continuity},   
    pdfauthor={Quang-Tuan Dang},   
    colorlinks=true,  
   linkcolor=purple,         
    citecolor=blue,        
    filecolor=green,      
    urlcolor=blue}          



\textwidth=360pt
 \textheight=615pt
 \parindent=8mm
 \evensidemargin=0pt
 \oddsidemargin=0pt
\frenchspacing
 \usepackage{fullpage}



\usepackage{colonequals}



\usepackage{enumitem}
% \usepackage{refcheck}
\renewcommand{\baselinestretch}{1.05}
\numberwithin{equation}{section}

%\usepackage{showkeys}


\newtheorem{thmA}{Theorem}
\renewcommand{\thethmA}{\Alph{thmA}}
\newtheorem{corA}[thmA]{Corollary}
%\renewcommand{\thecorA}{\Alph{corA}}

% theorems with special labels
\newtheorem{thmspec}{\relax}

\newtheorem{theorem}{Theorem}[section]
\newtheorem{proposition}[theorem]{Proposition}
\newtheorem{corollary}[theorem]{Corollary}
\newtheorem{lemma}[theorem]{Lemma}
\theoremstyle{definition}
\newtheorem{definition}[theorem]{Definition}
\newtheorem{remark}[theorem]{Remark}
\newtheorem{remarks}[theorem]{Remarks}
\newtheorem{example}[theorem]{Example}
\newtheorem{exercise}[theorem]{Exercise}
\newtheorem{question}[theorem]{Question}
\newtheorem{conjecture}[theorem]{Conjecture}
\newtheorem{problem}[theorem]{Problem}

\newcommand{\cali}[1]{\mathscr{#1}}

\newcommand{\GL}{{\rm GL}}
\newcommand{\Sp}{{\rm Sp}}
\renewcommand{\sp}{{\rm \mathfrak{sp}}}
\newcommand{\SO}{{\rm SO}}
\newcommand{\SU}{{\rm SU}}
\newcommand{\U}{{\rm U}}
\newcommand{\End}{{\rm End}}
\newcommand{\Vol}{{\rm Vol}}
\newcommand{\Hol}{{\rm Hol}}
\newcommand{\SL}{{\rm SL}}
\newcommand{\MA}{{\rm MA}}
\newcommand{\NMA}{{\rm NMA}}
\renewcommand{\sl}{{\rm\mathfrak{sl}}}
\newcommand{\Aut}{{\rm Aut}}
\newcommand{\lov}{{\rm lov}}
\newcommand{\supp}{{\rm Supp}}
\newcommand{\wi}{\widetilde}
\newcommand{\maxep}{\mathop{\mathrm{\rm  max_\epsilon}}\nolimits}
\newcommand{\maxunit}{\rm  max_1}
\newcommand{\const}{\mathop{\mathrm{const}}\nolimits}
\newcommand{\ev}{{\rm ev}}
\newcommand{\ov}{\widehat}
\newcommand{\Int}{{\rm Int}}

\newcommand{\diff}{\text{\normalfont d}}

\newcommand{\Leb}{\mathop{\mathrm{Leb}}\nolimits}
\newcommand{\Lip}{\mathop{\mathrm{Lip}}\nolimits}
\newcommand{\LLip}{{\mathop{\mathrm{\widetilde{Lip}}}\nolimits}}

 \newcommand{\dist}{\mathop{\mathrm{dist}}\nolimits}
\renewcommand{\Re}{\mathop{\mathrm{Re}}\nolimits}
\renewcommand{\Im}{\mathop{\mathrm{Im}}\nolimits}
\newcommand{\vol}{\mathop{\mathrm{vol}}}

\newcommand{\eq}{\mathrm{eq} }
%\newcommand{\ddc}{{dd^c}}
\newcommand{\ddc}{dd^c}
\newcommand{\dc}{d^c}
%\newcommand{\d}{\text{\normalfont d}}
%\newcommand{\Ta}{\text{\normalfont T}}
\newcommand{\Fix}{\text{\normalfont Fix}}
\newcommand{\Reg}{\text{\normalfont Reg}}
\newcommand{\Sing}{\text{\normalfont Sing}}






\newcommand{\mult}{{\rm mult}}
\newcommand{\Ssupp}{{\rm supp^*}}
\newcommand{\loc}{{loc}}
%\newcommand{\ddc}{dd^c}
%\newcommand{\dc}{d^c}
\newcommand{\diag}{{\rm diag}}
\newcommand{\prim}{{\rm prim}}
\newcommand{\otextbf}{{\rm \mathbf{0}}}
%
%
%
\def\d{\operatorname{d}}
\def\Ker{\operatorname{Ker}}
\def\mB{\mathcal{B}}
\def\mD{\mathcal{D}}
\def\mE{\mathcal{E}}
\def\mO{\mathcal{O}}
\def\mL{\mathcal{L}}
\def\mP{\mathcal{P}}
\def\mQ{\mathcal{Q}}
\def\mV{\mathcal{V}}
\def\mW{\mathcal{W}}
\def\mF{\mathcal{F}}
%
%
%
\DeclareMathOperator{\ch}{ch}
\DeclareMathOperator{\ric}{Ric}
\DeclareMathOperator{\spec}{Spec}

\newcommand{\Ta}{\text{\normalfont T}}
\newcommand{\dbar}{\widehat\partial}
\newcommand{\ddbar}{\partial\widehat\partial}
\renewcommand{\GL}{{\rm GL}}
\newcommand{\PSH}{{\rm PSH}}
\newcommand{\PB}{{\rm PB}}
\newcommand{\PC}{{\rm PC}}
\newcommand{\CN}{{\rm CN}}
\newcommand{\MM}{{\rm M}}
\newcommand{\ind}{{\bf 1}}
\newcommand{\id}{{\rm id}}
\newcommand{\Id}{{\rm Id}}
\newcommand{\pr}{{\rm pr}}

\newcommand{\ord}{{\rm ord}}
\newcommand{\exph}{{\rm exph}}
\newcommand{\capacity}{\mathop{\mathrm{Cap}}\nolimits}
%\newcommand{\capacity}{{\rm cap}}

\newcommand{\BT}{{\rm BT}}
\newcommand{\BTK}{{\rm BTK}}
\newcommand{\sing}{{\rm sing}}
\newcommand{\codim}{{\rm codim\ \!}}
\newcommand{\sign}{{\rm sign\ \!}}
\newcommand{\BB}{{\rm B}}
\newcommand{\HR}{{\rm HR}}
%\newcommand{\chac}{{\rm \bf 1}}
\newcommand{\chac}{{\rm \mathbf{1}}}
\newcommand{\Ac}{\cali{A}}
\newcommand{\Bc}{\cali{B}}
\newcommand{\Cc}{\cali{C}}
\newcommand{\Dc}{\cali{D}}
\newcommand{\Ec}{\cali{E}}
\newcommand{\Fc}{\cali{F}}
\newcommand{\Gc}{\cali{G}}
\newcommand{\Hc}{\cali{H}}
\newcommand{\K}{\cali{K}}
\newcommand{\Lc}{\cali{L}}
\newcommand{\Oc}{\cali{O}}
\newcommand{\Pc}{\cali{P}}
\newcommand{\Qc}{\cali{Q}}
\newcommand{\Rc}{\cali{R}}
\newcommand{\Uc}{\cali{U}}
\newcommand{\Vc}{\cali{V}}
\newcommand{\Sc}{\cali{S}}
\newcommand{\Tc}{\cali{T}}



\newcommand{\capK}{\text{Cap}}
%\newcommand{\span}{\rm span}


%\newcommand{\Z}{\mathbb{Z}}
\newcommand{\FS}{{\rm FS}}
\newcommand{\B}{\mathbb{B}}
\newcommand{\C}{\mathbb{C}}
\newcommand{\D}{\mathbb{D}}
\newcommand{\N}{\mathbb{N}}
\newcommand{\Z}{\mathbb{Z}}
\newcommand{\R}{\mathbb{R}}
\renewcommand\P{\mathbb{P}}

\renewcommand{\S}{\mathbb{S}}

\newcommand{\e}{\textbf{\textit{e}}}



%\newcommand{\K}{{\cal K}}
\newcommand{\E}{{\mathbb{E}}}
\renewcommand{\H}{\mathbb{ H}}

\newcommand{\swedge}{\widehat\wedge}
\newcommand{\transposee}[1]{{\vphantom{#1}}^{\mathit t}{#1}}
\newcommand{\spin}{$\text{spin}^c$ }
\newcommand{\norm}[1]{\lVert#1\rVert}
\newcommand{\abs}[1]{\lvert#1\rvert}

%------------------------------Boldsymbol-------------------------------------
\newcommand{\boldsym}[1]{\boldsymbol{#1}}
\newcommand\bb{\boldsym{b}}
\newcommand\bE{\boldsym{E}}
\newcommand\bg{\boldsym{g}}
\newcommand\bI{\boldsym{I}}
\newcommand\bJ{\boldsym{J}}
\newcommand\bk{\boldsym{k}}
\newcommand\bm{\boldsym{m}}
\newcommand\br{\boldsym{r}}
\newcommand\bt{\boldsym{t}}


\newcommand*{\doi}[1]{\href{http://dx.doi.org/#1}{doi: #1}}
\setcounter{tocdepth}{1}
%\setcounter{secnumdepth}{3}
%opening


\begin{document}
		\title[Modulus of continuity]{Modulus of continuity of Monge--Amp\`ere potentials in big cohomology classes}
	\author{Quang-Tuan Dang, Hoang-Son Do, Hoang Hiep Pham}
	
		
			\address{Yau Mathematical Sciences Center, Tsinghua University, Beijing 100084, China}
		\email{dangqt@mail.tsinghua.edu.cn $\&$ dangquangtuan10@gmail.com}
  \address{Vietnam Academy of Science and Technology, Institute of Mathematics, 18 Hoang Quoc Viet road, Cau Giay, Hanoi, Vietnam}
  \email{dhson@math.ac.vn}
  \address{Institute for Artificial Intelligence, University of Engineering and
Technology, Vietnam National University, Hanoi, Vietnam}
 \email{phhiepvn@gmail.com}
   % \address{Vietnam Academy of Science and Technology, Institute of Mathematics, 18 Hoang Quoc Viet	road, Cau Giay, Hanoi, Vietnam}
   % \email{phhiepvn@gmail.com}
	\date{\today}
	\keywords{Complex Monge-Amp\`ere equations; H\"older measure; Capacities}
%	\thanks{The work is partially supported by the ANR project PARAPLUI}
	\subjclass[2020]{32U15, 32W20, 32Q15}
	





\begin{abstract} In this paper, we prove a uniform estimate for the modulus of continuity of solutions to the degenerate complex Monge--Amp\`ere equation in big cohomology classes. This improves the previous work of Di Nezza--Lu and the first author.
 % ~\cite{dang2021continuity}.
\end{abstract}



\maketitle

  % \tableofcontents
\section{Introduction}




The study of complex Monge-Amp\`ere equations on compact K\"ahler manifolds has attracted considerable interest since Yau’s resolution of the Calabi conjecture. In connection with the Minimal Model Program, the search for singular K\"ahler--Einstein metrics
leads to the study of degenerate complex Monge–Amp\`ere equations; see~\cite{eyssidieux2009singular,berman2019kahler,guedj2017degenerate} and references therein.

Guedj and Zeriahi~\cite{guedj2007weighted} developed the first step of the study of the non-pluripolar Monge--Amp\`ere measure and solved degenerate complex Monge--Amp\`ere equations with rather general measures on the right-hand side. Their approach was later extended to the setting of big cohomology classes by Boucksom, Eyssidieux, Guedj, and Zeriahi~\cite{boucksom2010monge}. Since then, when the right-hand side belongs to $L^p$ for some $p>1$, the H\"older continuity of solutions to degenerate complex Monge--Amp\`ere equations has been intensively studied by many authors ~\cite{kolodziej1998complex,kolodziej2008holder,
hiep2010holder,
eyssidieux2011viscosity,boucksom2010monge,demailly2014holder, dinh2014characterization, 
dang2021continuity,dinh2022complex}. 
The modulus of continuity of the solution plays a crucial role since it is closely related to the geometric properties of the corresponding K\"ahler metrics, such as the uniform bounds for diameter of metrics and Gromov--Hausdorff convergence; see~\cite{Fu-Guo-Song2020-geometric,guo2022-local,guo2021modulus,guo2022-diameter, Guo2024-diameter2,guedj2025-diameter,vu2024continuity,do2023log,Guedj-To-25-green,vu24-diameter,nguyen-vu24-diameter}. 
 

The primary goal of this paper is to study the modulus of continuity of solutions to complex Monge--Ampère equations when the right-hand side is not integrable.
To state our result, we first introduce some notation and terminology. Let $X$ be a compact K\"ahler manifold of dimension $n$ equipped with a K\"ahler form $\omega_X$. We let $d$, $d^c$ denote the real differential operators on $X$ defined by $d:=\partial+\Bar{\partial}$, $d^c:=\frac{i}{2}(\Bar{\partial}-\partial)$ so that $\ddc=i\partial\Bar{\partial}$.
Fix a closed real smooth (1,1)-form $\theta$ representing a big cohomology class. Let $\PSH(X,\theta)$ denote the set of all $\theta$-psh functions. 
We say that the cohomology class $\{\theta\}$ is {\em big} if the set $\PSH(X,\theta-\varepsilon\omega_X)\neq \varnothing$ for some $\varepsilon>0$.
Let $\textrm{Amp}(\theta)$ denote the {\em ample locus} of $\theta$ which is roughly speaking the largest Zariski open subset where $\{\theta\}$ locally behaves like
a K\"ahler class. 

We are interested in studying the regularity of solutions to the complex Monge--Amp\`ere equation type $u\in\PSH(X,\theta)$ satisfying \begin{equation}\label{eq: dcmae}
   (\theta+\ddc u)^n=\nu,\quad\sup_X u=0, 
\end{equation}
where
$\nu$ is a positive measure on $X$ that puts no mass on pluripolar subsets and satisfies the compatibility condition $\nu(X)=\Vol(\theta)$, $u$ is the unknown $\theta$-function, and the left-hand side of~\eqref{eq: dcmae} denotes the non-pluripolar Monge--Amp\`ere product~\cite{bedford1987fine,boucksom2010monge}. 
% The existence and uniqueness of a solution $u$ with finite energy were shown in~\cite{boucksom2010monge}. When $\nu$ has density in $L^{1+\varepsilon}$, by~\cite{boucksom2010monge,demailly2014holder} such a solution $u$ has minimal singularities and H\"older continuous in some Zariski open subset. 

Let $\mathcal{C}$ denote the set of $\omega_X$-psh functions $u$ normalized by $\int_X u\omega_X^n=0$. This is a
convex compact set in the $L^1(\omega_X^n)$ topology. Define the following distance on $\mathcal{C}$ 
\[\dist(u,v)\coloneqq \|u-v\|_{L^1}, \]
where the $L^1$-norm is with respect to the measure $\omega_X^n$. 
We say that a measure $\mu$ is {\em H\"older continuous}
	 with H\"older constant $B$ and H\"older exponent $0<\beta\leq 1$
	 with respect to $\dist_{L^1}$ on $\mathcal{C}$ if for any $u,v\in\mathcal{C}$
     \[\int_X|u-v|\omega_X^n\leq B\|u-v\|_{L^1}^\beta. \]
Our main result is the following:
\begin{theorem}\label{main} Let $(X,\omega_X)$ be a compact K\"ahler manifold of dimension $n$ and $\theta$ a smooth closed (1,1) form whose cohomology class is big.
	 Let $\mu$ be a H\"older continuous measure 
	 with H\"older constant $B$ and H\"older exponent $0<\beta\leq 1$
	 with respect to $\dist_{L^1}$ on $\mathcal{C}$.
	  Assume that $\psi$ is a quasi-psh function
	 on $X$ satisfying $\int_Xe^{-\psi}d\mu=\Vol(\theta)$ and $\omega_X+a_0dd^c\psi\geq 0$,
	 with $a_0>0$. Suppose that  
	 $u\in\mathcal{E} (X, \theta)$ is a solution to the following equation \begin{equation} (\theta+\ddc u)^n=e^{-\psi}\mu,\quad\sup_X u=0.\end{equation}
	  Then for each $U\Subset{\rm Amp}(\theta)\backslash\{\psi=-\infty\}$, there exists a continuous function $F_U:\R_+\rightarrow\R_+$
	  with $F(0)=0$ such that
	  $$|u(z_1)-u(z_2)|\leq F_U(\dist (z_1, z_2)),$$
	  for every $z_1, z_2\in U$. Moreover, the choice of $F_U$ depends only on  $X$, $U$, $\omega_X$, $n$, $\theta$, $a_0$, $B$, $\beta$, $\sup_U(-\psi)$ and
	  an upper bound function for  $H(a)=\int_Xe^{{2(V_{\theta}-u)}/{a}}d\mu$.
\end{theorem}
Here we notice that $e^{{2(V_{\theta}-u)}/{a}} \in L^1(\mu)$ for every $a>0$ by Lemma~\ref{lem integ phi-u}. In the case $\theta=\omega_X$,
we have the following corollary:
% {\color{red} Do we prove that $u$ is H\"older continuous ? not only continuous}
\begin{corollary}\label{coro: main}Let $(X,\omega_X)$ be a compact K\"ahler manifold of dimension $n$.
	Let $\mu$ be a H\"older continuous measure 
	with H\"older constant $B$ and H\"older exponent $0<\beta\leq 1$
	with respect to $dist_{L^1}$ on $\mathcal{C}$.
	Assume that $\psi$ is a quasi-psh function
	on $X$ satisfying $\int_Xe^{-\psi}d\mu=\Vol(\theta)$, $\omega_X+a_0dd^c\psi\geq 0$ and \begin{equation}\label{eq: H}
	    \int_Xh(-\psi)e^{-\psi}\mu\leq C_0,
	\end{equation}
	where $a_0, C_0>0$ are constants and $h: (0, \infty)\rightarrow (0, \infty)$ is an increasing, 
	concave function with $h(\infty)=\infty$. Suppose that  
	$u\in\mathcal{E} (X, \omega_X)$ is a solution to the following equation \begin{equation}\label{eq: cmae_kahler} (\omega_X+\ddc u)^n=e^{-\psi}\mu,\quad\sup_X u=0.\end{equation}
	Then for each $U\Subset X\backslash\{\psi=-\infty\}$, there exists a continuous function $F_U:\R_+\rightarrow\R_+$
	with $F(0)=0$ such that
	$$|u(z_1)-u(z_2)|\leq F_U(\dist (z_1, z_2)),$$
	for every $z_1, z_2\in U$.  Moreover, the choice of $F_U$ depends only on $X$, $U$, $\omega_X$, $n$, $\theta$, $a_0$, $B$, $\beta$, $\sup_U(-\psi)$, $C_0$,
	 and $h$.
\end{corollary}
Let us emphasize that the continuity of the solution $u$ to equation~\eqref{eq: cmae_kahler} was established by Di Nezza and Lu~\cite[Theorem 3.1]{di2017complex}. The contribution of the corollary above is to establish the equicontinuity of the family of solutions corresponding to pairs $(\mu,\psi)$ which satisfy the condition~\eqref{eq: H}, given $h$ and $C_0$. We also note that for each pair $(\mu,\psi)$ with $e^{-\psi}\in L^1(X,\mu)$, there always exists a pair $(h,C_0)$ such that the condition~\eqref{eq: H} holds; cf. Remark~\ref{rmk: hC}. 
%Furthermore, we show that the constant in Di Nezza and Lu’s estimate is independent of any upper bound on $\int_X e^{-2u/a}d\mu$.
% In~\cite[Theorem 3.1]{di2017complex}, Di Nezza and Lu proved that the solution $u$ to the equation~\eqref{eq: cmae_kahler} is continuous, whereas our theorem establishes its modulus of continuity. Furthermore, we show that the constant in Di Nezza and Lu's estimate is independent of an upper bound on $\int_X e^{-2u/a}d\mu$.




\subsection*{Acknowledgment.}  
This work was done while the authors were visiting the Vietnam Institute for Advanced Study in Mathematics (VIASM), and they would like to thank VIASM for its hospitality and support.
% \subsection*{Ethics declarations} The authors declare no conflict of interest.
\section{Preliminaries} 
% We recall some essential materials that can be found in the literature~\cite{demaillycomplex,demailly1994regularization}.
Throughout the paper, we let $X$ denote a compact K\"ahler manifold of dimension $n$, equipped with a K\"ahler form $\omega_X$. 
We denote by \(dV_{\omega_X}:={\omega_X^n} \)
the volume form associated with $\omega_X$. For any measure $\mu$ on $X$, we write $L^1(\mu)$ for $L^1(X,d\mu)$.
% We will work in local coordinates $(z_1,\ldots,z_n)$ and write a tensor in terms of its components in such a coordinate system. 
% \subsection{K\"ahler geometry and quasi-psh functions} 

\subsection{Quasi-plurisubharmonic functions}
%  A Hermitian metric $h$
%  on $X$
% is a positive definite Hermitian $\mathcal{C}^\infty$ form on $TX$; in local coordinates can be written by $h=\sum_{j,k}g_{j\Bar{k}}dz_j\otimes d\Bar{z}_k$.
% We associate it with $\omega_X$ a  K\"ahler metric given in local coordinates by
% \[\omega_X=-\textrm{Im}h=\frac{i}{2}\sum_{j,k=1}^ng_{k\Bar{j}}dz_k\wedge d\Bar{z}_{j}.\]
% We normalize by $\int_X\omega_X^n=1$. Let $D$ be the Levi-Civita (Chern) connection associated with $g$ (they coincide in the K\"ahler case). Let also $\Theta=\Theta(D)=D^2$ the curvature tensor. It is known (see~\cite{demaillycomplex}) that
% $i\Theta$ is a section of the buundle $\mathcal{C}^\infty_{1,1}(X, \textrm{Herm}(TX,TX))$. From analytic point of view, $i\Theta$ is locally a matrix with $(1,1)$-forms as entries. In terms of local coordinates, we can write
% \[i\Theta=i\sum_{j,k,l,m}c_{jklm}dz_j\wedge d\Bar{z}_k\otimes \left( \frac{\partial}{\partial z_l}\right)^* \otimes\frac{\partial}{\partial z_m}. \]
% We can identify the curvature tensor to a Hermitian form 
% \[i\Theta(\xi \otimes v):=\sum_{j,k,l,m=1}^nc_{jklm} \xi_j\Bar{\xi}_k v_l\Bar{v}_m\] and $c_{kjml}=\Bar{c}_{jklm}$, $c_{jklm} = c_{lkjm} = c_{jmlk} = c_{lmjk}$ by symmetry. The K\"ahlerness implies that these coefficients are exactly the one of the bisectional (Levi-Civita) curvature tensor of $g$ which is locally defined by
% \[c_{jklm}=R_{j\Bar{k}l\Bar{m}}:=-\frac{\partial^2 g_{l\Bar{m}}}{\partial z_j\partial\Bar{z}_j}+\sum_{p,q=1}^ng^{p\Bar{q}}\frac{\partial g_{l\Bar{q}}}{\partial z_j}\frac{\partial g_{p\Bar{m}}}{\partial \Bar{z}_k}.\]
% We say that the bisectional curvature is bounded from below by $-A$ for $A>0$ if for any $z\in X$ and any vectors $a,b\in T_z X\setminus \{0\}$ we have
% \[\sum_{j,k,l,m=1}^nR_{j\Bar{k}l\Bar{m}}(z)a_j\Bar{a}_kb_l\Bar{b}_m\geq -A\|a\|^2_g\cdot\|b\|^2_g. \]
% The existence of such a bound does not depend on the choice of local coordinates.
% \medskip

Recall that an upper semi-continuous function $ \varphi:X \rightarrow\mathbb{R}\cup\{-\infty\} $
is called {\it quasi-plurisubharmonic} ({\it quasi-psh} for short) if it is locally the sum of a smooth and a plurisubharmonic (psh for short) function. 

We say that $\varphi$ is {\it $\theta$-plurisubharmonic}  ({\it $\theta$-psh} for short) if it is quasi-psh, and $$\theta_\varphi:=\theta+\ddc \varphi\geq 0$$ in the sense of currents, where ${\rm d}=
\partial+\Bar{\partial}$ and ${\rm d}^c=\frac{i}{2\pi}(\Bar{\partial}-\partial)$ so that $\ddc=\frac{i}{\pi}\partial\Bar{\partial}$. 
% By the $\ddc$-lemma any closed  positive $(1,1)$-current  $T$ cohomologous to $\theta$ can be written as $T=\theta+\ddc\varphi$ for some $\theta$-psh function $\varphi$ which is furthermore unique up to an additive constant. 
We let $\PSH(X,\theta)$ denote the set of all $\theta$-psh functions, which are not identically $-\infty$. This set is endowed with 
the weak topology, which coincides  with 
 the $L^1$-topology. By Hartogs' lemma, $\varphi\mapsto\sup_X\varphi$ is continuous in this weak topology, 
 % Since the set of  closed positive currents  in a fixed  cohomology class  is compact (in the weak topology), 
it follows that the set of $\varphi\in\PSH(X,\theta)$, with $\sup_X\varphi=0$ is compact. 
We refer the reader to~\cite{demaillycomplex,guedj2017degenerate} for more properties of $\theta$-psh functions.   


The cohomology class $\{\theta\}$ is said to be {\em big} if the set $\PSH(X,\theta-\varepsilon\omega_X)$ is not empty for some $\varepsilon>0$.	
	We now assume that $\{\theta\}$ is big unless otherwise specified.
By Demailly's regularization theorem \cite{demailly1992regularization}, we can find a closed positive $(1,1)$-current $T_0\in \{\theta\}$ such that $$T_0=\theta+\ddc\Psi_0\geq \varepsilon_0\omega_X,$$ for some $\varepsilon_0>0$, where $\Psi_0$ is a quasi-psh function with \emph{analytic singularities}, i.e., locally %and its global potential $\f$ are said to have \emph{analytic singularities} if there exists $c>0$ such that 
	$$	\Psi_0=c\log\left[\sum_{j=1}^{N}|f_j|^2\right]+g,
	$$
	%locally on $X$,	
	where the $f_j$'s are holomorphic functions and $g$ is bounded. Such a current $T_0$ is then smooth on a Zariski open subset $X\setminus\{\Psi_0=-\infty\}$. We thus define the {\em ample locus} $\textrm{Amp}(\theta)$ of $\theta$ as the largest such Zariski open subset (which exists by the Noetherian property of closed analytic subsets; cf.~\cite{boucksom2004divisorial}).


Given $\varphi,\psi\in \PSH(X,\theta)$, we say that $\varphi$ is {\it less singular} than  $\psi$, and denote by $\psi\preceq\varphi$, if  there exists a constant $C$ such that $\psi\leq \varphi+C$ on $X$. We say that $\varphi,\psi$ have the {\em same singularity type}, and denote by $\varphi\simeq\psi$ if $\varphi\preceq\psi$ and $\psi\preceq \varphi$. 
% This defines an equivalence relation on
% $\PSH(X,\theta)$, whose equivalence classes are the singularity types $[\varphi]$. 
There is a
natural least singular potential in $\PSH(X,\theta)$ given by
\begin{align*}
	V_{\theta}:=\sup\{\varphi\in \PSH(X,\theta): \varphi\leq 0\}.
	\end{align*}
A function $\varphi$ is said to have {\em minimal singularities} if it has the same singularity type as $V_\theta$. In particular, $V_\theta=0$ if $\theta$ is semi-positive.
We see that $V_\theta$ is locally bounded in the ample locus.

 % Let $x_0\in X$. Fixing a holomorphic chart $x_0\in V_{x_0}\subset X$,  the {\em Lelong number} $\nu(\varphi,x_0)$ of a quasi-psh function $\varphi$ at  $x_0\in X$ is defined as follows:
	% 	\begin{align*}
	% 	\nu(\varphi,x_0):=\sup\{\gamma\geq 0: \varphi(z)\leq \gamma\log\|z-x_0\|+O(1), \; \text{on}\; V_{x_0}\}.  
	% 	\end{align*}
 % Remark that this definition does not depend on the choice of local holomorphic charts.
\subsection{Demailly's regularization} 


   We consider the exponential mapping $ \exp_x: T_x X\ni \zeta\to X$ as follows: $\exp_x(\zeta)=\gamma(1)$ where $\gamma:[0,1]\to X$ is the geodesic starting from $x=\gamma(0)$ with the initial velocity  $\gamma'(0)=\zeta$. In the Euclidean space $\mathbb{C}^n$, the exponential map $\exp_x(\zeta)=x+\zeta$. 
 % Unfortunately, The exponential Riemannian map is, unfortunately, not holomorphic on a general complex manifold. To overcome this difficulty Demailly~\cite{demailly1994regularization} introduced the new mapping $\exph$ as follows.
% \begin{definition}
%     A $\mathcal{C}^\infty$ map $\exph:TX\to X$ is defined by
%     \begin{itemize}
%         \item For every $x\in X$, $\exph_x(0)=x$ and $d_\zeta\exph_x(0)={\rm Id}_{T_x X}$;
%         \item  For every $x\in X$ the map $\zeta\mapsto\exph_x\zeta$ has a holomorphic Taylor expansion at $\zeta=0$ of the form
%         \[\left| \exph_z(\zeta)_m-z_m-\zeta_m-\frac{1}{2}\sum_{j,k,\ell }R_{j\Bar{k}\ell\Bar{m}}(\Bar{z}_k+\frac{1}{3}\Bar{\zeta}_k)\zeta_j\zeta_\ell \right|\leq O(|\zeta|^2(|z|+|\zeta|)^2),\quad |\zeta|<r \]
%         for $r>0$ sufficiently small, where we have chosen local normal coordinates $(z_1,\ldots,z_n)$ on $X$ centered at $x$ and holomorphic normal coordinates $(\zeta_j)$ on the fibers of $TX$ near $x$ 
% with respect to $\omega_X$, and $R_{j\Bar{k}\ell\Bar{m}}$ are coefficients of bisectional curvature tensor of $\omega_X$ in these local coordinates.
%     \end{itemize}    
% \end{definition}
% We should mention here that the "holomorphic" exponential mapping above is the holomorphic counterpart of the exponential mapping $\exp$; it is defined as the holomorphic part of the Taylor expansion of $\zeta\mapsto\exp_x(\zeta)$ thanks to E. Borel's theorem. Note also that the mapping is unique, up to a $\mathcal{C}^\infty$ functions flat along the zero section of $TX$ which is identified with $X$. As mentioned by Demailly, the $\exph$ mapping is still not holomorphic (since the Taylor expansion at $0$ may be divergent). Nevertheless, it satisfies the so-called quasi-holomorphic properties that will be enough for our purpose. 

% \medskip

% Let $\exph\colon TX\to X$, $ \exph_x: T_x X\ni \zeta\to X$  be  the formal holomorphic part
% of the Taylor expansion of the exponential map of the Chern connection on $TX$ associated with the metric $\omega_X$.

 Following~\cite{demailly1982estimations,demailly1994regularization}, for a quasi-psh function $u$,
we define its {\em $\delta$-regularization} $\rho_\delta u=\Psi(u)(z,\delta)$ where
\begin{equation}\label{phie}
\Psi(u)(z,w)=\int_{\zeta\in T_{z}X}
u\big(\exp_z(w\zeta)\big)\chi\big({|\zeta|^2_{\omega_X }}\big)\,dV_{\omega_X}(\zeta),\ \delta>0,
\end{equation}
where $\chi: \mathbb R_{+}\rightarrow\mathbb R_{+}$ is defined by
\begin{center}
$\chi(t)=\begin{cases}\frac {\eta}{(1-t)^2}\exp(\frac 1{t-1})&\ {\rm if}\ 0\leq t\leq 1,\\0&\
	{\rm if}\ t>1,\end{cases}$
\end{center}
with a suitable constant $\eta$ such that
$\int_{\mathbb C^n}\chi(\Vert z\Vert^2)\,dV(z)=1$, where $dV$ denotes the Lebesgue measure on $\mathbb{C}^n$. The $\delta$-regularization $\rho_\delta u$ can be written by
\[\rho_\delta u(z)=\frac{1}{\delta^{2n}}\int_{\zeta\in T_{z}X}
u\big(\exp_z(\zeta)\big)\chi\bigg(\frac{|\zeta|^2_{\omega_X }}{\delta^2}\bigg)\,dV_{\omega_X}(\zeta).\]
Intuitively, this corresponds to the familiar convolution with smoothing kernels. Actually, in the case of $\mathbb{C}^n$ endowed with the Euclidean metric, this is exactly the smoothing convolution; see~
\cite[Remark 4.6]{demailly1994regularization}. 

  
% The proof itself was originated from [Dem94, Proposition 3.8] and [BD12, Lemma 1.12]. We restate the result here 
The following lemma is a combination of \cite[Theorem 4.1]{demailly1994regularization} and \cite[Lemma 1.12]{berman2012regularity}.
\begin{lemma}\label{lem kiselman}
	Let $u$ be a $\theta$-psh function and define the Kiselman-Legendre transform with level $c$ by
	\begin{equation}\label{kisleg}
	\Phi_{c,\delta}\coloneqq\inf _{0< t\leq \delta }\Big[\rho_tu(z)+K (t^2-\delta^2)+K (t-\delta)
	-c\log\Big(\frac t{\delta}\Big)\Big].
	\end{equation}	
	Then for  some positive constants $0<\delta_0<1$ and $K>1$  depending on the curvature tensor of $\omega_X$ on $X$,
	 the function $\rho_tu(z)+Kt^2$ is increasing in $t\in (0, \delta_0]$ and
	one has the following estimate for the complex Hessian:
	\begin{equation}\label{hessest}
	\theta+dd^c \Phi_{c,\delta}\geq -(Ac+K^2\delta)\,\omega_X,
	\end{equation}
	where $A$ is a lower bound of the negative part of the bisectional curvature of $\omega_X$.
\end{lemma}
% To the authors' knowledge, the above lemma was already proved in~\cite[Lemma 4.1]{kolodziej2019stability}. However, for readers' convenience, we briefly sketch the proof which claims that the lemma still holds without the boundedness of $u$. By the explanation in \cite[page 624]{demailly2014holder} the proof holds for the exponential function $\exp_x$ instead of its holomorphic part $\exph_x$.
%  {\color{red} For the proof of this lemma, we require at least that $u$ is locally bounded on an open dense subset $U\subset X$ (in this case $\lambda(z,\delta)=0$). We can show that the estimate holds for any  $x\in U$, then extend it to everywhere. Unless we need to restart as follows.
% \[\theta+dd^c \Phi_{c,\delta}\geq -(A\min(c,\lambda(z,\delta))+2K\delta)\,\omega_X. \]
% }
  \begin{proof} 
  % The lemma was already proved in
 The proof is identical to that of~\cite[Lemma 4.1]{kolodziej2019stability}, which is still valid without the boundedness of $u$.
 \end{proof}
%  We borrow notation from~\cite{demailly1994regularization}. 
%  Given $\vartheta$, $\eta$ tangent vectors in the coordinates $z\in \mathbb{C}^n$, $w\in\mathbb{C}$, we simply denote by $[\vartheta]^2$ and  $[\eta]^2$ the (1,1) vectors $\vartheta\wedge\Bar{\vartheta}$ and $\eta\wedge\Bar{\eta}$. Note that if we write $\theta=\sum_j\vartheta_j\frac{\partial}{\partial z_j}$ then $d|z|^2[\vartheta]^2=\sum_{j=1}^n|\vartheta_j|^2$, similarly $|dw|^2[\eta]^2=|\eta|^2$.
% At a point $(x,0)$ we choose a local holomorphic coordinates $(z,w)$ centered at $(x,0)$. We argue the same as in~\cite[Theorem 4.1]{demailly1994regularization}. Let $g$ be a local smooth potential of $\theta$. We see that $\Psi(g+u)=g+\Psi(u)+O(|w|^2)$. We can thus work on a small coordinate open set $U\subset X$.

% The estimate (4.3) in ~\cite{demailly1994regularization} implies that
% \begin{align*}
%     \theta &+\ddc\Psi(z,w)[\vartheta,\eta]^2\geq\\
%     &|w|^2\int_{\mathbb{C}^n}-\chi_1(\zeta\sum_{j,k,l,m}\frac{\partial^2 u}{\partial z_m\partial\Bar{z}_l})(\exph_z(w\zeta))\left( c_{jklm}+\frac{1}{|w|^2}\delta_{jm}\delta_{kl}\right)\tau_j\Bar{\tau}_k dV(z)\\
%     &\quad -K'(|\vartheta||\eta|+|\eta|^2)
% \end{align*}
% where $c_{jklm}$ denote the coefficients of the curvature w.r.t $\omega_X$, $\chi_1$ is the primitive $-\chi_1(t)=\int_t^{+\infty}\chi(s)ds$ and $\tau=\vartheta+\zeta\eta+O(|w|)$. Let us mention that the estimate above is obtained by ~\cite[(3.11)]{demailly1994regularization} (applied with $N=2$). 
% We let $A=\sup_{|a|,|b|\leq 1}-c_{jklm}a_j\Bar{a}_kb_l\Bar{b}_m$. The above estimate can be written as
% \begin{align*}
%     &\theta+\ddc\Psi(z,w)[\vartheta,\eta]^2\geq\\
%     &-A \left( |w|^2\int_{\mathbb{C}^n}-\chi_1(\zeta\sum_{j,k,l,m}\frac{\partial^2 u}{\partial z_m\partial\Bar{z}_l})(\exph_z(w\zeta)) dV(z)\right)|\vartheta|^2-K''(|\vartheta||\eta|+|\eta|^2)\\
%     &\geq -A\lambda_U(z,w)|\vartheta|^2-K''(|\vartheta||\eta|+|\eta|^2).
% \end{align*} where $\lambda_U(z,w)$ is defined by
% \[\lambda_U(z,w)=|w|^2\int_{\mathbb{C}^n}-\chi_1(|\zeta|^2)\frac{\partial^2 u}{\partial z_k\partial\Bar{z}_l}(z+w\zeta) dV(\zeta). \] We remark that $K',K''$ are $O(1)$. We expand the function $\exph_z$ in the Taylor series up to order $N=2$ and apply~\cite[Prop. 2.9 (ii)]{demailly1994regularization}. The coefficients of the Taylor expansion of the exponential map at 0 only depend on  the curvature of $\omega_X$ and its subsequents covariant derivatives at $x$, thus they are uniformly bounded because $X$ is compact. It follows that
% \begin{align*}
%     &\theta+\ddc\Psi(z,w)[\vartheta,\eta]^2\geq\\
%     &\geq -A\lambda_U(z,w)|\vartheta|^2-K''(|\vartheta||\eta|+|\eta|^2)\\
%     &-A\lambda_U(z,w)|\vartheta|^2-C|w||\vartheta|^2-K'''(|\vartheta||\eta|+|\eta|^2)
% \end{align*} Recall that
% \[\lambda(z,w):=\frac{\partial(\Psi(u)(z,t)+Kt^2)}{\partial\log t}\big|_{t=|w|}=t\partial_t(\Psi+Kt^2
% )\xrightarrow{t\to 0^+} \nu(u,x).\]
% From \cite[(4.5)]{demailly1994regularization} we have $|\lambda(z,w)-\lambda_U(z,w)|=O(|w|^2)$. Thus
 % \[ \theta+\ddc\Psi(z,w)[\vartheta,\eta]^2\geq-(A\lambda(z,w)+K'''|w|^2)|\vartheta|^2-C|w||\vartheta|^2-K'''(|\vartheta||\eta|+|\eta|^2).\]
% Putting $\Tilde{K}=K'''+C$ we conclude that
% As above, we choose local normal coordinates $(z,\zeta)$ with respect to $\omega_X$ on the tangent bundle of $X$ at some point $(x,0)$. 
% According the estimates in~\cite[Section 8.1]{demailly1982estimations},
% the result of~\cite[Theorem 4.1]{demailly1994regularization} is still valid in the K\"ahler setting, so that there are constant $\delta>0$, $\Tilde{K}>1$ such that for $(z,w)\in X\times\mathbb C$
% \begin{equation*}
%    \theta+\ddc \Psi(z,w)\geq -A\lambda(z,w) |dz|^2-\Tilde{K}(|w|^2+|w|)|dz|^2 -\Tilde{K}|dz||dw|-\Tilde{K}|dw|^2,
% \end{equation*}
% where $A=\sup_{|a|\leq 1,|b|\leq 1}\{-c_{jk\ell m}a_j\Bar{a}_kb_\ell\Bar{b}_m\}$, $c_{jk\ell m}$ are the coefficients of the curvature tensor of $(T_X,\omega_X)$ for $K>0$ to be chosen hereafter and 
% \[ \lambda(z,w)\coloneqq \frac{\partial(\Psi(u)(z,t)+\tilde Kt^2)}{\partial\log t}\Big|_{t=|w|}=t\partial_t(\Psi+\tilde Kt^2
%  )\xrightarrow{t\to 0^+} \nu(u,x).\]
% We have that
% \[\Tilde{K}|dz||dw|\leq \Tilde{K}^2|w||dz|^2 +\frac{1}{4|w|}|dw|^2.\]
% Fix $c>0$. We set 
% \[U_{c,\delta}(z,w):=\Psi(u)(z,w)+ K(|w|^2-\delta^2)+ K(|w|-\delta)-c\log\frac{|w|}{\delta}. \]
% and
% \[\Phi_{c,\delta}:=\inf_{0<|w|<\delta}U_{c,\delta}(z,w). \]
% Since $c\log|w|$ is pluriharmonic for $|w|>0$,  simple computations give
% \begin{align*}
%    \theta+\ddc U_{c,\delta}&=\theta+\ddc \Psi(z,w)+
%    K(1+1/4|w|)|dw|^2 \\
%    &\geq -\left( A\lambda(z,|w|) +\Tilde{K}|w|^2+ (\Tilde{K}^2+\Tilde{K})|w| \right)|dz|^2 \\
%    &\quad+\left( K\bigg(1+\frac{1}{4|w|}\bigg)-\Tilde{K} -\frac{1}{4|w|}\right)|dw|^2\\
%    &\geq -\left( A\lambda(z,|w|) +K|w| \right)|dz|^2 
% \end{align*} with $K$ being chosen so that $K\geq \Tilde{K}^2+2\Tilde K+1$. 
% By definition of $\lambda$ and $U_{c,\delta}$ we have that
% \[ \frac{\partial U_{c,\delta}}{\partial
% \log t}=\lambda(z,t)+ Kt-c\xrightarrow{t\to 0} -c.\]
% The function $t\mapsto U_{c,\delta}(z,t)$ is strictly increasing near 0 then this function attains a minimum value in $(0,\delta]$ at some point $t_0=t_0(z)$. If $t_0=\delta$ then $\lambda(z,t_0)+K\delta\leq c$  in particular $\lambda(z,t_0)<c$. Otherwise $0<t_0<\delta$ the derivative vanishes then $\lambda(z,t_0)+Kt_0=c$ so $\lambda(z,t_0)<c$. From both cases, by Kiselman' principle we  infer that 
% \[ \theta+\ddc\Phi_{c,\delta}\geq -(Ac+K\delta)\omega_X\] in the ample locus, hence everywhere on $X$. 

\begin{remark}
   The above lemma is essential for our study of the modulus of continuity of solutions to complex Monge--Ampère equations. We follow the same strategy as in the proof of~\cite[Theorem D]{demailly2014holder}. 
    % However, this proof had a gap when they applied the incorrect lemma in~\cite[Lemma 1.12]{berman2012regularity}. 
    However, the function $\Phi_{c,\delta}$ differs from the one used there by the addition of the term $K(t - \delta)$, as introduced in~\cite[Remark 4.7]{demailly1994regularization}. This extra term is needed to handle the mixed term $|dz||dw|$ that arises in the estimates; cf.~\cite{kolodziej2019stability}.
\end{remark}
\begin{lemma}\label{Jen}
	% Assume that $u$ is a continuous $\omega_X$-psh function
	% normalized by  $\min_{X}u =1,\ \max_{X}u\leq B$ for some fixed constant $B$. 
 Let $u\in\PSH(X,\theta)$.
 If
	$\rho_{\delta} u$ is the regularization of
	$u$ defined as in (\ref{phie}) then for $\delta$ small enough we have
	$$
	\int_{X} \frac{\rho_{\delta} u-u}{\delta ^2 }\omega_X^n \leq C\|u\|_{L^1(X,dV_X)} ,
	$$
	where $C$ only depends on $n$, $X$, $\omega_X$.
\end{lemma}
\begin{proof}
    % We consider a double covering $U_i\Subset V_i$, $i=1,...,N$. In each local chart $V_i$, we write $\theta=\ddc h_i$ for some potential $h_i$. For sufficiently small $\delta$ we have the inclusion
    % \[\{ \exph_z(\zeta): z\in U_i, \|\zeta\|<\delta\}\subset V_i. \]
    % As explained, one can make approximation with $\psi_i=u_i+h_i$ (where $u_i$ is the restriction of $u$ on $U_i$)
    % whenever the base point varies in $U_i$. Note that it follows from \cite[Remark 4.6]{demailly1994regularization} that  $\rho_\delta u(z)=u*\chi_\delta+O(\delta^2)$. 
    % It thus suffices to get control on the $L^1$-average of $(\psi_{i})_\delta-\psi_i$ on $U_i$. An application of Jensen's type formulae gives us the desired estimate; see~\cite[Lemma 4.3]{guedj2008holder}.
    % {\color{red} more details here!!!}
    The proof follows verbatim from that of~\cite[Lemma 2.3]{demailly2014holder}.
%     , so we briefly give the main idea here. Since $x\mapsto\zeta=\log_z x$ is the inverse of $\zeta\mapsto x =\exp_z(\zeta)$, so by definition
%     \begin{align*}
%         \rho_\delta u=\int_{x\in X} u(x)\rho\left(\frac{|\log_z x|^2_{\omega_X}}{\delta^2} \right)\frac{dV_{\omega_X}(\log_z x)}{\delta^{2n}}=\int_X u(x)K_\delta(z,x)
%     \end{align*} where we set $$K_\delta(z,x)=\rho\left(\frac{|\log_z x|^2_{\omega_X}}{\delta^2} \right)\frac{dV_{\omega_X}(\log_z x)}{\delta^{2n}}.$$
%     The latter is a semi-positive $(n,n)$-form on $X\times X$; it can be viewed as a kernel with
% compact support in a neighborhood of the diagonal of $X\times X$. We have that $\int_XK_\delta(z,x)=1$. Hence by the change of variable $(z,x)\mapsto (x,z)$
% \begin{align*}
%  \int_X   (\rho_\delta u(z)-u(z))dV_{\omega_X}&=\int_{X\times X}(u(x)-u(z))K_\delta(z,x)\delta dV_{\omega_X}(z)\\
%  &=\int_{X\times X}u(x)(K_\delta(z,x)\wedge dV_{\omega_X}(z)-K_\delta(x,z)\wedge dV_{\omega_X}(x)).
% \end{align*}
% Thanks to Fubini's theorem and ~\cite[Lemma 2.4]{demailly2014holder}, the latter integral is equal
% to \begin{align*}
%     \int_X\int_{z\in B(x,\delta)} & u(x)(K_\delta(z,x)\wedge dV_{\omega_X}(z)-K_\delta(x,z)\wedge dV_{\omega_X}(x))\\
%     & \leq \int_X\int_{z\in B(x,\delta)} |u(x)|C\delta^{2-2n}dV_{\omega_X}(z)\wedge dV_{\omega_X}(x)\\
%     &\leq C'\delta^2\int_X |u(x)|dV_{\omega_X}(x).
% \end{align*}
% If $\mu$ is H\"older continuous measure then
% \[\int_X(\rho_\delta u-u)d\mu\leq C_0B\delta^{2\beta}\|u\|^\beta_{L^1(X,\omega_X)}. \]
\end{proof}
% \subsection{Approximations of quasi-psh functions}
% \begin{theorem}[\cite{demailly1992regularization,demailly2015cohomology}]\label{the Dem appro}
% 	Let $\varphi$ be a $\theta$-psh function on $X$. There exists a decreasing sequence of quasi-psh functions $(\varphi_m)$ such that
% 	\begin{enumerate}
% 			\item $\varphi_m$ have the same
% 		singularities as  $1/2m$ times a logarithm of a sum of squares of holomorphic functions,
% 		\item $(\varphi_m)$ converges  to $\varphi$ as $m\to+\infty$, { pointwise and in $L^1$};
% 		\item $\theta+dd^c\varphi_m\geq -\varepsilon_m\omega_X$, where $\varepsilon_m>0$ decreases to 0 as $m\to+\infty$;
% 		\item $\int_Xe^{2m(\varphi_m-\varphi)}dV<+\infty$;
% 		\item $\varphi_m$ is smooth outside the analytic subset $E_{1/m}(\varphi)$.
% 	\end{enumerate}
% \end{theorem}

% Since the cohomology class $\{\theta\}$ is big, there exists a negative 
% $\theta$-psh function $\Psi_0$ with analytic
% singularities such that $\theta +dd^c\Psi_0\geq\varepsilon_0\omega_X$
% for some $\varepsilon_0>0$. 
% It follows from~\cite{boucksom2004divisorial} that we can choose $\psi$ so that ${\rm Amp}(\theta)=X\setminus {\rm Sing}(\Psi_0)$. By subtracting a positive constant, we can assume $\Psi_0\leq V_\theta$.

% Fix $0<c_1\leq \frac{\varepsilon_0}{4A}$ 
% and $0<\delta_1<\min\{\delta_0, \sqrt{\frac{\varepsilon_0}{4K}}\}$, where $\delta_0$ is defined as in Lemma \ref{lem kiselman}. For every $0<c<c_1$ and $0<\delta<\delta_1$,
% we define
% \begin{equation}\label{eq: constantB0}
% B_0=B_0(c, \delta)=\frac{Ac+K^2\delta}{\varepsilon_0}.
% \end{equation}
% Then, we have $0<B_0<1/2$.  We set
% \begin{equation}\label{eq: u c delta}
% u_{c, \delta}=B_0\Psi_0+(1-B_0)\Phi_{c, \delta}.
% \end{equation}
% where $\Phi_{c, \delta}$ is defined by \eqref{kisleg}. It follows from Lemma~\ref{lem kiselman} that $\theta+\dc u_{c,\delta} \geq \varepsilon_0B_0^2\omega_X$ hence $u_{c,\delta}$ is $\theta$-psh. From now on, we set $\Omega:={\rm Amp}(\theta)\setminus\{\psi=-\infty\}$.


% {\color{red} Need a proof here.}

% By Proposition~\ref{prop DN14 moderate} and \cite[Theorem 3.1]{guedj2021quasi1}, one has:
% \begin{corollary}\label{cor bounded}
% 	Let $\mu$ be a positive measure on $X$ of mass $\int_X(\theta+dd^cV_{\theta})^n$.
% 	 If $\mu$ is H\"older continuous then there exists a unique $u\in \PSH(X, \theta)$ with
% 	 minimal singularities such that $(\theta+dd^cu)^n=\mu$ and $\sup_Xu=0$. Moreover, $V_{\theta}-C\leq u\leq V_{\theta}$ for some $C>0$.
% \end{corollary}
\subsection{Non-pluripolar product}


Let $\theta^1,\ldots,\theta^n$ be closed smooth real (1,1) form representing big cohomology classes, and $\varphi_j\in\PSH(X,\theta^j)$. Following the construction of Bedford--Taylor~\cite{bedford1987fine}, it has been shown in~\cite{boucksom2010monge} that for each $k\in\mathbb{N}$,
\[ \mathbf{1}_{\cap_j\{\varphi_j>V_{\theta^j}-k\}}(\theta^1+\ddc{\max(\varphi_1,V_{\theta^1}-k)})\wedge\cdots\wedge (\theta^n+\ddc{\max(\varphi_n, V_{\theta^n}-k)})\] is well-defined as a  Borel positive measure with finite total mass. The sequence of these measures is non-decreasing in $k$ and it converges weakly to the so-called {\em Monge--Amp\`ere product}, denoted by \[ (\theta^1+\ddc{\varphi_1})\wedge \cdots\wedge (\theta^n+\ddc{\varphi_n}),\] 
which does not charge pluripolar sets by definition. In particular, $\theta^1=\cdots=\theta^n=\theta$ and $\varphi_1=\cdots=\varphi_n$ we obtain the non-pluripolar
Monge--Amp\`ere measure of $\varphi$, denoted by $(\theta+\ddc\varphi)^n$ or simply by $\theta_\varphi^n$.

The {\em volume} of a big cohomology class $\{\theta\}$ is given by the total mass of the non-pluripolar Monge--Ampère
measure of $V_\theta$, i.e.,  $${\rm \Vol}(\theta):=\int_{X}\theta_{V_\theta}^n.$$
  We say that $\varphi\in\PSH(X,\theta)$ has {\it full Monge--Amp\`ere mass} if $\int_X\theta_{\varphi}^n=\Vol(\theta)$. We let \begin{align*}
      \mathcal{E}(X,\theta):=\left\{\varphi\in\PSH(X,\theta):\int_X\theta_{\varphi}^n=\Vol(\theta) \right\}
  \end{align*}
  denote the set of $\theta$-psh functions with full Monge--Amp\`ere mass.
  Note that $\theta$-psh functions with minimal singularities have full  Monge--Amp\`ere mass (see \cite[Theorem 1.16]{boucksom2010monge} for more details), but the converse is not true. 

\smallskip 
% We recall the result of Dinh-Nguyen in K\"ahler case.
% \begin{theorem}[{\cite{dinh2014characterization}}]
% 	Let $\mu$ be a positive measure on $X$ of mass $\int_X\omega_X^n$. If $\mu$ is H\"older continuous
% 	with H\"older exponent $0<\beta\leq 1$ with respect to ${\rm dist}_{L^1}$ on $\mathcal{C}$ then
% 	 the equation
% 	\begin{equation}\label{eq the DN14}
% 	(\omega_X+dd^cu)^n=\mu,
% 	\end{equation}
% 	admits a H\"older continuous solution with H\"older exponent $\alpha$ for all $0<\alpha<\frac{2\beta}{n+1}$.
% 	Conversely, if \eqref{eq the DN14} admits a H\"older continuous solution then $\mu$ is H\"older continuous.
% \end{theorem}



Given a measurable function $f:X\to\mathbb{R}$, we define the {\em $\theta$-psh envelope} of $h$ by
\begin{equation*}
P_\theta(f):=(\sup\{u\in\PSH(X,\theta): u\leq f\;\text{on}\, X \})^*,
\end{equation*} where the star means we take the upper semi-continuous regularization. %We next define 
%\[ P_{\theta}(g,h):=(\sup\{u\in\PSH(X,\theta): u\leq \min (g,h)\;\text{on}\, X \})^*.\] 

Given a $\theta$-psh function $\phi$, 
Ross and Witt-Nystr\"om~\cite{ross2014analytic} introduced the ``rooftop envelope'' as follows
\[P_\theta[\phi](f)= \left( \lim_{C\to +\infty}P_\theta(\min(\phi+C,f))\right)^*. \]
When $f=0$ we simply write $P_\theta[\phi]$. 
A function $\phi\in\PSH(X,\theta)$ is called a {\em model potential} if $$\int_X\theta_\phi^n>0\quad\text{and}\quad\phi=P_\theta[\phi].$$
% \subsection{Capacity}

We recall the notion of Monge--Ampère capacity introduced in~\cite{di2017complex,di2015generalized}.
For $\phi \in \PSH(X,\theta)$ and a Borel subset $E \subset X$, the {\em relative Monge–Ampère capacity} is defined by
    \begin{equation*}
        \capacity_{\phi}(E):=\sup\left\{\int_E\theta_u^n: u\in\PSH(X,\theta),\;\phi-1\leq u\leq \phi \right\}.
    \end{equation*}
This definition recovers the Monge--Ampère capacity used in~\cite{boucksom2010monge} when $\phi = V_\theta$. In that case, we denote it simply by $\capacity_\theta$.

\begin{definition}
A family of positive measures $\{\mu_i \}_{i\in I}$ on $X$ is said to be uniformly absolutely
continuous with respect to $\phi$-capacity if, for every $\varepsilon>0$, there exists $\delta>0$ such that for each Borel subset $E\subset X$ satisfying $\capacity_\phi(E)<\delta$ the inequality $\mu_i(E)<\varepsilon$ holds for all $i\in I$. We denote this by $\mu_i<<\capacity_\phi$. In particular, all such measures must
vanish on pluripolar sets.
\end{definition}
% \begin{proposition}
%     There exists a continuous function $g:\mathbb
%     R^+\to\mathbb R^+$ with $g(0)=0$ such that for any Borel set $E\subset X$, 
%     \[m(E) \leq g(\textrm{Cap}_\phi(E))\; \forall\, m\in\mathcal{S}. \]
% \end{proposition}

\subsection{H\"older continuous measures}
\begin{definition}\label{def: holder}
	A positive measure $\mu$ on $X$ is said to be $\PSH(X, \theta)$-H\"older continuous
	(or, for short, H\"older continuous) if there exist $B>0$ and $0<\beta\leq 1$ such that
	\begin{equation*}
	    \int_{X}(u-v)d\mu\leq B \left(\int_X|u-v|\omega_X^n\right)^{\beta},
	\end{equation*}
for every $u, v\in\mathcal{C}:=\{w\in \PSH(X, \theta): \int_X w\omega_X^n=0 \}.$ In this case, we say that $B$ is
the H\"older constant and
$\beta$ is the H\"older exponent of $\mu$ with respect to ${\rm dist}_{L^1}$ on $\mathcal{C}$. We let $\mathcal{M}(B,\beta)$ denote the set of such measures.
\end{definition}
\begin{lemma}[{\cite[Lemma 3.3]{dinh2014characterization}}]\label{lem DN14 def holder}
		Assume $\mu$ is a H\"older continuous measure 
		with H\"older constant B and H\"older exponent $0<\beta\leq 1$
		 with respect to $\dist_{L^1}$ on $\mathcal{C}$. Then there exists $C>0$ depending only on $B$ such that
		\begin{equation*}
		    \|u-v\|_{L^1(\mu)}\leq C\max\big\{\|u-v\|_{L^1(\omega_X^n)}, \|u-v\|_{L^1(\omega_X^n)}^{\beta} \big\},
		\end{equation*}
	for all $u, v\in \PSH(X, \theta)$. In particular, if a family $\mathcal{F}$  of $\theta$-psh functions
	is a relatively compact subset of $L^1(\omega_X^n)$ then there exists $C_{\mathcal{F}}>0$ such that
	\begin{equation*}
	    \|u-v\|_{L^1(\mu)}\leq C_{\mathcal{F}} \|u-v\|_{L^1(\omega_X^n)}^{\beta},
	\end{equation*}
	for all $u, v\in\mathcal{F}$.
\end{lemma}


\begin{proposition}\label{prop exp integrable}
	Let $\mu$ be a H\"older continuous measure on $X$
	with H\"oler constant $B$ and H\"older exponent $0<\beta\leq 1$ with respect to ${\rm dist}_{L^1}$ on $\mathcal{C}$. Assume that
	$u, v$ are $\theta$-psh functions satisfying
	\begin{equation*}
	    \int_Xe^{m(u- v)}\omega_X^n<C_0,
	\end{equation*}
for positive constants $m$, $C_0$. Then, for every $0<\gamma<\beta$, there exists $C>0$ depending only on $C_0, 
m, B,\beta$ and $\gamma$ such that
\begin{equation*}
    \int_Xe^{\gamma m(u-v)}d\mu<C.
\end{equation*}
\end{proposition}
\begin{proof} We adapt the proof of~\cite[Proposition 4.4]{dinh2014characterization}.
	Denote $w=v-u$. We have
	\begin{equation}\label{eq0 exp integ}
		\omega_X^n(\{w<-M\})\leq\int_Xe^{-m(w+M)}\omega_X^n\leq C_0e^{-mM},
	\end{equation}
for every $M>0$. Denote 
	$$w_M=\max\{w, -M \}=\max\{v, u-M\}-u.$$
We have $w_M-w=\max\{v, u-M\}-v$ is the difference of two $\theta$-psh functions. 
By Lemma \ref{lem DN14 def holder}, there exists $C_1>0$ depending only on $B$ such that
\begin{equation}\label{eq1 exp integ}
\int_X(w_M-w)d\mu\leq C_1\max\left\{\|w_M-w\|_{L^1(\omega_X^n)}, \|w_M-w\|_{L^1(\omega_X^n)}^{\beta} \right\}.
\end{equation}
Moreover, we also have
\begin{equation}\label{eq2 exp integ}
\mu(\{w<-M-1\})\leq\int_X(w_M-w)d\mu,
\end{equation}
and by \eqref{eq0 exp integ}, 
\begin{equation}\label{eq3 exp integ}
\begin{split}
    \int_X(w_M-w)\omega_X^n&=\int_0^{\infty}\omega_X^n(\{w<-M-t\})dt\\
    &\leq C_0\int_0^{\infty}e^{-m(M+t)}dt
=\dfrac{C_0e^{-mM}}{m}.
\end{split}
\end{equation}
Combining \eqref{eq1 exp integ}, \eqref{eq2 exp integ} and \eqref{eq3 exp integ}, we get

$$	\mu(\{w<-M-1\})\leq C_1\max\left\{\dfrac{C_0e^{-mM}}{m}, \left(\dfrac{C_0e^{-mM}}{m}\right)^{\beta}\right\}
	\leq C_2\cdot e^{-\beta mM},$$
where $C_2=C_1\max\{\frac{C_0}{m}, (\frac{C_0}{m})^{\beta}\}$. Then, for every $0<\gamma<\beta$, we obtain
\begin{align*}
	\int_Xe^{-\gamma m w}d\mu-\mu (X)
	&=\gamma m\int_0^{\infty}\mu(w<-t)e^{\gamma m t}dt\\
	&\leq 	\gamma m C_2e^{\beta m}\int_0^{\infty}e^{(-\beta+\gamma)m t}dt\\
	&=\frac{C_2\gamma e^{\beta m}}{\beta-\gamma}.
\end{align*}
\end{proof}



\begin{lemma}\label{lem integ phi-u} Let $\mu\in\mathcal{M}(B,\beta)$ be a H\"older continuous measure on $X$
	and $u\in\mathcal{E}(X, \theta)$.
	Then
		$$\int_Xe^{m(V_{\theta}-u)}d\mu<+\infty,$$
for every $m>0$.
 \end{lemma} 
 \begin{proof} 
     It follows from~\cite[Proposition 2.10]{dang2021continuity} that for any $b>0$, $P_{\omega_X}(b(u-V_\theta))\in\mathcal{E}(X,\omega_X)$. We remark here that $u-V_\theta$ is well-defined outside a pluripolar set. 
     In particular, for $0<\gamma<\beta$, we have
     \[\int_X e^{m\gamma^{-1}(V_\theta-u)}\omega_X^n\leq \int_X e^{-m\gamma^{-1}P_{\omega_X}(u-V_\theta)}\omega_X^n<+\infty, \]
     as follows from Skoda's integrability theorem; cf.~\cite[Theorem 2.50]{guedj2017degenerate}. We, therefore, apply Proposition~\ref{prop exp integrable} to conclude.
 \end{proof}



 
\begin{proposition}\label{the exp cap}
	Let $\phi\in \PSH(X, \theta)$ be a model potential with $\int_X\theta_{\phi}^n>\varrho>0$. 
    % Assume that $\mu$ is a positive measure on $X$ which is H\"older continuous	with H\"older constant $B$ and  H\"older exponent $0<\beta\leq 1$ with respect to $\dist_{L^1}$ on $\mathcal{C}$. 
    Let $\mu\in\mathcal{M}(B,\beta)$.
    Then, there exist constants
	$0<\gamma<1$ and $C>0$ depending only on $X, \omega_X, \theta, B$ and $\beta$ such that
	\begin{equation}\label{eq0 the exp cap}
		\mu (E)\leq C\,
		\exp\left(-\gamma\left(\dfrac{\varrho}{\capacity_{\phi}(E)}\right)^{1/n} \right),
	\end{equation}
	for all Borel sets $E\subset X$. 
\end{proposition}
\begin{proof} The case when $\theta=\omega_X$ and $\phi=0$ was shown in \cite[Propositions 2.4, 4.4]{dinh2014characterization}. For the relative case,
   the proof relies on the arguments in~\cite[Proposition 4.30]{darvas2018monotonicity}, so we sketch it here. 
   According to~\cite[Section 4A2]{darvas2018monotonicity}, for any Borel set $E$, we define the global $\phi$-extremal function of $(E,\theta,\phi)$ by
   \[V_{E,\phi}=\sup\{ u\in\PSH(X,\theta,\phi): u\leq \phi \;\text{on}\, E\}. \]
   % We define also the relative Alexander–Taylor capacity of $E$ by \[T_\phi(E)\coloneqq\exp(-M_\phi(E)),\quad\text{where}\, M_\phi(E)=\sup_X V^*_{E,\phi}. \] 
 Set $M_\phi(E)=\sup_X V^*_{E,\phi}$, we may assume $M_\phi(E)\geq 1$.  Thanks to~\cite[Lemma 4.9]{darvas2018monotonicity}, we have the Alexander--Taylor type inequality
 \[ \left(\frac{\varrho}{\capacity_{\phi}(E)}\right)^{1/n} \leq M_\phi(E),\]
 hence,
 \[ \exp(-M_\phi(E))\leq \exp\left(-\left(\frac{\varrho}{\capacity_{\phi}(E)}\right)^{1/n} \right). \] 
 % Noticing that we can assume $M_\phi(E)\geq 1$.
  As follows from~\cite[Proposition 4.4]{dinh2014characterization}, $\mu$ is weakly moderate, i.e. there are constants $\alpha=\alpha(\beta)>0$ and $C=C(B,\beta)>0$ such that $\int_X \exp(-\alpha \varphi)d\mu\leq C$ for every $\varphi\in\mathcal 
   C$. We apply this to $V_{E,\phi}^*-\int_X V_{E,\phi}^*\omega_X^n$ to obtain 
   \[\int_X\exp(-\alpha V_{E,\phi}^*)d\mu\leq C\cdot \exp\bigg(-\alpha\int_XV_{E,\phi}^*\omega_X^n\bigg). \]
   Since $V_{E,\phi}^*\leq 0$ on $E$ a.e. and $\int_X V_{E,\phi}^*{\omega_X}^n\geq M_\phi(E)-C_{\omega_X}$ (see, e.g.,~\cite[Proposition 8.5]{guedj2017degenerate}) we have
   \[ \mu(E)\leq C\cdot\exp\left(-\gamma\left(\frac{\varrho}{\capacity_{\phi}(E)}\right)^{1/n} \right),\] for $\gamma>0$. 
\end{proof}

\begin{proposition} \label{prop: unif-capa}
% Let $\mu$ be a positive measure on $X$ which is H\"older continuous with H\"older constant $B$ and  H\"older exponent $0<\beta\leq 1$ with respect to $\dist_{L^1}$ on $\mathcal{C}$.
Fix $a_0>0$, $ C_0>0$ and $h: (0, \infty)\rightarrow (0, \infty)$ is an increasing concave function with $h(\infty)=\infty$.
Let $\mathcal{N}=\mathcal{N}(B,\beta, a_0,C_0,h)$ be the set of probability measures $\nu$ on $X$ such that $\nu=e^{-\psi}\mu$ with $\mu\in\mathcal{M}(B,\beta)$, $\psi\in\PSH(X, a_0\omega_X)$ and  \begin{equation}\label{eq: assumption H}
    \int_Xh(-\psi)e^{-\psi}{\rm d}\mu\leq C_0.
\end{equation}  
% where $\psi$ is a quasi-psh function on $X$ 
Then there exists a continuous function $g\colon [0,\infty)\to [0,\infty)$ with $g(0)=0$ such that for all Borel sets $E$,
\[\nu(E)\leq g(\capacity_{\omega_X}(E)) \quad\text{for all}\,\nu\in\mathcal{N}.\]
In particular, the family of measures $(\nu)_{\nu\in\mathcal{N}}$ is uniformly absolutely
continuous with respect to capacity.
\end{proposition}

\begin{proof}
    For any Borel subset $E\subset X$, if $\capacity_{\omega_X}(E)=0$, then $E$ is a pluripolar set, and consequently $\nu(E)=0$. In this case, the desired inequality holds trivially.
We may therefore assume $\capacity_{\omega_X}(E)>0$. We have for $k>0$
    \begin{align*}
        \int_E e^{-\psi}d\mu&=\int_{E\cap\{\psi\geq -k \} } e^{-\psi}d\mu+\int_{E\cap \{\psi<-k \}}e^{-\psi}d\mu\\
        &\leq e^k\mu(E)+ \frac{1}{h(k)}\int_{E\cap \{\psi<-k \}}h(-\psi)e^{-\psi}d\mu\\
        &\leq C(e^k\capacity_{\omega_X}(E)+{h(k)}^{-1}),
    \end{align*}
    where the last inequality follows from Proposition~\ref{the exp cap}. 
    If $\capacity_{\omega_X}(E)=0$ then $E$ is a pluripolar set, so $\nu(E)=0$, the desired inequality is trivial.
 Otherwise, taking $k\coloneqq\log \frac{1}{\sqrt{\capacity_{\omega_X}(E)}}>0$ we get that $e^{-\psi}\mu(E)\leq g(\capacity_{\omega_X}(E))$ where
    \[ g(t)\coloneqq C(t^{1/2}+h(\log (t^{-1/2}))^{-1}).\]
     Otherwise, since $\frac{1}{\sqrt{\capacity_{\omega_X}(E)}}\leq 1$, the choice of $g$ ensures that $\nu(E)\leq C\leq g(\capacity_{\omega_X}(E))$.
\end{proof}
\begin{remark}\label{rmk: hC}
    For each non-pluripolar measure $e^{-\psi}\mu$ with $\int_Xe^{-\psi}d\mu<\infty$, there always exists a concave increasing function $h\colon \mathbb{R}^+
    \to\mathbb R^+$ and $C>0$ such that $\int_X h(-\psi)e^{-\psi}d\mu\leq C$ as follows from~\cite[Lemma 3.3]{boucksom2010monge}. 
\end{remark}
% \section{H\"older continuous measures}
We end this section with some examples of H\"older continuous measures.
\begin{example}
   Dinh and Nguy\^en~\cite{dinh2014characterization} gave a characterization of H"older continuous measures in the sense of Definition~\ref{def: holder}. Specifically, they proved that a measure $\mu$ is H\"older continuous if and only if it can be expressed as the Monge--Amp`ere measure of a H\"older continuous $\omega_X$-psh function, i.e., $\mu = (\omega_X + \ddc\varphi)^n$ for some $\varphi\in\PSH(X,\omega_X)$ that is H\"older continuous. As a consequence, one obtains explicit examples of such measures; several are listed below.
    \begin{itemize}
        \item $\mu=f\nu$ where $\nu$ is a H\"older continuous measure and $f\in L^p(\nu)$ for some $p>1$; cf.~\cite{kolodziej2008holder,demailly2014holder}. 
        \item If there are constants $C>0$ and $\alpha>0$ such that $\mu(B(z,r))\leq Cr^{2n-2+\alpha}$ for $B(z,r)$ denoting the ball centered at $z$ with radius $r>0$, then $\mu$ is H\"older continuous; cf.~\cite{hiep2010holder,demailly2014holder}.
        \item If $\mu$ is a positive Radon measure compactly supported on a immersed $\mathcal{C}^3$ real submanifold $K$ of $X$ of real codimension $d>0$, then $\mu$ is H\"older continuous; cf.~\cite{Vu2018complex}.
    \end{itemize}
\end{example}
\begin{example} We construct an example of a measure with \emph{radial singularities}, inspired by~\cite[Section 1.3]{di2021finite} and~\cite[Section 5]{guedj2025-diameter}.
Let $X=\mathbb C\mathbb P^n$ be a complex projective manifold of dimension $n$, equipped with
the Fubini--Study metric $\omega_X=\omega_{\rm FS}$. We assume that ${\varphi}$ is a $\omega_{\rm FS}$-psh function on $X$ which has a radial singularity at a point $p$, i.e., it is invariant under the group $U(n,\mathbb C)$ in the neighborhood of $p$. Choosing a local chart biholomorphic to the unit ball $B\subset \mathbb{C}^n$, with $p$ corresponding to the origin. Locally, the function $\varphi$ can be written as $\varphi=u-\frac{1}{2}\log[1+\|z\|^2]$ for some psh function $u$.

We therefore consider a psh function $u\coloneqq\chi\circ L$ defined near the origin in $\mathbb C^n$, where $L(z)\coloneqq\log \|z\|$ and $\chi\colon\mathbb R^-\to \mathbb R^- $ is a convex increasing function. A computation shows that
\[ (\ddc u)^n=\frac{c_n (\chi'\circ L)^{n-1}\chi''\circ L}{\|z\|^{2n}}\omega_X^n=\colon e^{-\psi_\chi}\omega_X^n. \]
To simplify later estimates, we assume that $\chi(t)$ does not go to zero too fast at infinity $-\infty$, that is, $\chi'(t),\chi''(t)\geq e^{Ct}$ near $t=-\infty$. Under this hypothesis, one finds $\psi_\chi\sim \log\|z\|$ near the origin. Hence,
the function $h$ satisfies $\int_X h(-\psi_\chi)e^{-\psi_\chi}\omega_X^n<\infty $ if and only if
\[\int_{-\infty}^{-1} h(-t)\chi'(t)^{n-1}\chi''(t)dt<\infty.\]
Below, we provide several families of functions $\chi$ for which the condition holds with $h(s)=s^\varepsilon$, where $0<\varepsilon\ll 1$.
\begin{enumerate}
    % \item Consider $\chi_a(t)=-(-t)^{-a}$ for some $a>0$. We observe that
    % \[\int_{-\infty}^{-1}h(-t)\chi'(t)^{n-1}\chi''(t)dt\lesssim\int_{-\infty}^{-1} h(-t)\cdot (-t)^{-na-n-1}dt<\infty \]
    % for $h(s)=s^p$ with $p<n(a+1)$.
    \item Consider $\chi_a(t)=(-t)^a$ for $a\in (0,1)$.  We have
    \[\int_{-\infty}^{-1}h(-t)\chi'(t)^{n-1}\chi''(t)dt\lesssim\int_{-\infty}^{-1} h(-t)\cdot (-t)^{-n(1-a)-1}dt<\infty, \]
    for $h(s)=s^\varepsilon$ with $0<\varepsilon<<1$ and $a<\frac{n-\varepsilon}{n}$.
    \item Consider $\chi_a(t)=-(\log (-t))^a$, where $a>0$. We have 
  \[\int_{-\infty}^{-1}h(-t)\chi'(t)^{n-1}\chi''(t)dt\lesssim\int_{-\infty}^{-1} h(-t)\cdot (-t)^{-n-1}(\log(-t))^{na-n}dt<\infty,\]
  for $h(s)=s^\varepsilon$ with $0<\varepsilon<<1$.
  \item $\chi(t)=-\log(\log(-t))$. We have
  \[\int_{-\infty}^{-1}h(-t)\chi'(t)^{n-1}\chi''(t)dt\lesssim\int_{-\infty}^{-1} h(-t)\cdot (-t)^{-n-1}(\log(-t))^{-n}dt<\infty,\]
   for $h(s)=s^\varepsilon$ with $0<\varepsilon<<1$.
\end{enumerate}
\end{example}

\section{Proof of the main theorem}
 We prove some lemmas which will be used to prove the main theorem.
% In this section, we prove Theorem \ref{main} by using  the Kiselman-Legendre transform.

% {\color{red} It's quite ambitious to get the above estimate since $V_\theta-u$ is not psh in general. However, the envelop technique could help. We first show that $P_\omega_X(b(u-V_\theta))$ is a $\omega_X$-psh function, i.e. not $-\infty$ for any $b>0$. Next, by monotonicity, we show that the latter function has full M-A mass. This yields $P_\omega_X(b(u-V_\theta))$ has zero Lelong numbers. The results follow from Skoda's integrability theorem. We can restate our lemma as follows. If $u\in\mathcal{E}(X,\theta)$ then
% % \[\int_Xe^{-P_\omega_X(m(u-V_\theta))}dV<+\infty.\]
% % See~\cite[Prop. 2.10]{Dan} for a proof. This leads us to modify Prop.~\ref{prop exp integrable} by setting $w=P_{\omega_X}(v-u)$ ?}

\begin{lemma}\label{lem t^n cap}
	Assume $u, v$ are negative $\theta$-psh functions such that
	\begin{itemize}
		\item $v\leq P_{\theta}[u]$;
		\item $v=\lambda v_1+(1-\lambda) v_2$, where  $v_1, v_2$ are $\theta$-psh functions
		with $v_2$ having model singularity type and $\lambda\in (0, 1)$. 
	\end{itemize}
Then, for every $s>0$ and $0\leq t\leq 1-\lambda$, we have
	$$t^n \capacity_{v_2}\{u<v-s- t\}\leq \int_{\{u<v-s\}}\theta_u^n.$$
\end{lemma}
% {\color{red} Why we need $v_2$ to be model type singularity; $|v_2-P_\theta[v_2]|\leq C$ ?}
\begin{proof} The proof follows closely that of \cite[Lemma 4.31]{darvas2018monotonicity}, which itself originates in \cite[Lemma 4.2]{boucksom2010monge}.
	 Let $\varphi$ be a $\theta$-psh function such that $v_2-1\leq\varphi\leq v_2$. For every
	 $s>0$ and $0\leq t\leq 1$, we have
	 \begin{align*}
	 	\{u<v-s-(1-\lambda) t\}\subset \{u< \lambda v_1+(1-t)(1-\lambda) v_2+(1-\lambda) t\varphi-s \}
	 	 \subset \{u<v-s\}.
	 \end{align*}
 Setting $\widehat{v}:=\lambda v_1+(1-t)(1-\lambda) v_2+(1-\lambda) t\varphi$, we have
 \begin{align*}
 	(1-\lambda)^nt^n\int_{\{u<v-s-(1-\lambda) t\}}\theta_{\varphi}^n
 	\leq \int_{\{u<v-s-(1-\lambda) t\}}\theta_{\widehat{v}}^n
 	\leq \int_{\{u<\widehat{v}-s\}}\theta_{\widehat{v}}^n,
 \end{align*}
for all $s>0$ and $0\leq t\leq 1$. Since $\widehat{v}\leq P_\theta[u]$
it follows from the comparison principle
 \cite[Lemma 2.3]{darvas2021log} that 
 \begin{align*}
   \int_{\{u<\widehat{v}-s\}}\theta_{\widehat{v}}^n  \leq \int_{\{u<\widehat{v}-s\}}\theta_{u}^n
 	\leq \int_{\{u<v-s\}}\theta_{u}^n.
 \end{align*}
 Since $\varphi$ was taken arbitrarily, it follows that
 	$$(1-\lambda)^nt^n \capacity_{v_2}\{u<v-s-(1-\lambda) t\}\leq \int_{\{u<v-s\}}\theta_u^n,$$
for all $s>0$ and $0\leq t\leq 1$. Substituting $t\mapsto (1-\lambda) t$, we obtain the desired inequality.
\end{proof}
% The following lemma goes back to De Giorgi.
\begin{lemma}\label{lem tg(t+s)}
	Let $g: \R^+\rightarrow\R^+$ be a non-increasing, right-continuous function such that, 
	\begin{align*}
		tg(t+s)\leq C(g(s))^{1+\alpha},
	\end{align*}
for every $s>0$ and $0\leq t\leq 1-\lambda$, where $C>0$, $\alpha> 0$ and $\lambda\in (0, 1)$ are given.
Assume that there exist $S_0>0$ and $0\leq t_0\leq 1-\lambda$ satisfying $g(S_0)^{\alpha}<\frac{t_0}{2C^{1/n}}$.
Then $g\left(S\right)=0$ for all $S\geq S_0+\frac{t_0}{1-2^{-\alpha}}$.
\end{lemma}
\begin{proof}
     We refer to \cite[pages 614–615]{eyssidieux2009singular} for a proof.
\end{proof}
\begin{lemma}\label{lem DL17}
	Let $u\in \mathcal{E} (X, \theta)$ such that $\theta_u^n=e^{-\psi}d\mu$, where $\mu$ is a H\"older continuous measure with H\"older constant $B$ and H\"older exponent $0<\beta\leq 1$ on $\mathcal{C}$
	and $\psi$ is a negative quasi-psh function. Suppose that $v$ is a $\theta$-psh function such that
	$v\leq V_{\theta}+a\psi$ for some $a>0$. Assume that $\int_X\theta_{V_{\theta}}^n>\varrho>0$, where $\varrho$ is a constant. Then,  there exists $C=C(n, X, \omega_X, \theta, \varrho, B, \beta)>0$ such that,
	for every $0<\lambda\leq  1/2$,
	\begin{align*}
		u\geq  \lambda v+(1-\lambda) V_{\theta}-C a\lambda \left(1+\log_+\int_Xe^{\frac{2(V_\theta-u)}{a\lambda}}d\mu\right)-2.
	\end{align*}
\end{lemma}
Here we note that $e^{\frac{2(V_{\theta}-u)}{a\lambda}} \in L^1(\mu)$ by Lemma~\ref{lem integ phi-u}. Denote by $\log_+(x)=\max(\log x,0)$.
\begin{proof}
	Denote $\widehat{v}\coloneqq \lambda v+(1-\lambda) V_{\theta}$.
	 We apply \ref{lem t^n cap} with $v_1=v$ and $v_2=V_\theta$ to obtain that for all $s>0$ and $0\leq t\leq 1-\lambda$,
	\begin{align*}
		t^n\capacity_{\theta}(u<\widehat{v}-s-t)\leq \int_{\{u<\widehat{v}-s\}}\theta_u^n
		=\int_{\{u<\widehat{v}-s\}}e^{-\psi}d\mu.
	\end{align*}
Since $v\leq V_{\theta}+a\psi$, we have $-\psi\leq \frac{V_\theta-\widehat v}{a\lambda}$.
It follows that
\begin{align*}
t^n\capacity_{\theta}(u<\widehat{v}-s-t)&\leq \int_{\{u<\widehat{v}-s\}}e^{\frac{V_{\theta}-\widehat{v}}{a\lambda}}d\mu\\
&\leq \int_{\{u<\widehat{v}-s\}}e^{\frac{V_{\theta}-u-s}{a\lambda}}d\mu.
\end{align*}
Then, by H\"older's inequality, we obtain
\begin{equation}\label{eq1 proof lem DL17}
t^n\capacity_{\theta}(u<\widehat{v}-s-t)\leq 
\left( \int_{X}e^{\frac{2(V_{\theta}-u-s)}{a\lambda}}d\mu\right)^{{1}/{2}}
\left( \int_{\{u<\widehat{v}-s\}}d\mu\right)^{{1}/{2}}.
\end{equation}
Since the measure $\mu$ is H\"older continuous, it follows from Proposition~\ref{the exp cap} that there exists $C_1>0$
depending on $X, \omega_X, \theta, \varrho, B$ and  $\beta$ such that
\begin{center}
	$\mu(K)\leq C_1^{2n} [\capacity_{\theta}(K)]^4,$
\end{center}
for every Borel set $K\subset X$, using that $\exp(e^{-1/x})=O(x^4)$ for all $x>0$. 
Then, by \eqref{eq1 proof lem DL17}, we have
\begin{equation}\label{eq2 proof lem DL17}
t^n\capacity_{\theta}(u<\widehat{v}-s-t)\leq C_1^n e^{\frac{-s}{a\lambda}}\left( \int_{X}e^{\frac{2(V_{\theta}-u)}{a\lambda}}d\mu\right)^{{1}/{2}}
[\capacity_{\theta}(u<\widehat{v}-s)]^2.
\end{equation}
Set $g(s)\coloneqq[\capacity_{\theta}(u<\widehat{v}-s)]^{1/n}$ and
 $C_2\coloneqq\left(\int_{X}e^{\frac{2(V_{\theta}-u)}{a\lambda}}d\mu\right)^{{1}/{2n}}$.
  By \eqref{eq2 proof lem DL17},
 we have
 \begin{equation}\label{eq3 proof lem DL17}
 	tg(t+s)\leq C_1C_2g(s)^2,
 \end{equation}
for every $s>0$ an $0\leq t\leq 1-\lambda$. Moreover, it follows from \eqref{eq1 proof lem DL17} (with $t=1-\lambda$) that
\begin{equation}\label{eq4 proof lem DL17}
	g(s+1-\lambda)\leq\dfrac{C_3}{1-\lambda} e^{\frac{-s}{a\lambda n}},
\end{equation}
for every $s>0$, where $C_3:=C_2(\mu(X))^{1/2n}$. Put
	$s_0=a\lambda n \log_+ (16C_1C_2C_3)$ so that
$$\frac{C_3}{1-\lambda} e^{\frac{-s_0}{a\lambda n}}<\frac{(1-\lambda)}{4C_1C_2}.$$
Then, by \eqref{eq4 proof lem DL17}, we have $g(s_0+1-\lambda)<\frac{(1-\lambda)}{4C_1C_2}$. 
Applying Lemma \ref{lem tg(t+s)} with $S=s_0+1-\lambda$ and $t_0=\frac{1-\lambda}{2}\in (0, 1-\lambda)$, we obtain
$g(s_0+2-2\lambda)=0$. Then
\begin{align*}
	u\geq\widehat{v}-s_0-2+2\lambda\geq \widehat{v}
	-C_4 a\lambda \left(1+\log_+\int_Xe^{\frac{2(V_\theta-u)}{a\lambda}}d\mu\right)-2,
\end{align*}
where $C_4>0$ is a constant depending only on $n, X, \omega_X, \theta, \varrho, B$ and $\beta$.
\end{proof}
\begin{lemma}\label{lem sup<cap}
		Let $u, v$ be negative $\theta$-psh functions such that
	\begin{itemize}
		\item $v\leq u+M$ for some $M\geq 1$;
		\item $v=\lambda v_1+(1-\lambda) v_2$, where  $v_1, v_2$ are $\theta$-psh functions, and $v_2$ has model singularity type and $\lambda\in (0, 1/2)$;
		\item there exist $C>0$ and $\alpha>0$ such that, for every $s>0$,
		\begin{equation}\label{eq0lemsup<cap}
			\int_{\{u<v-s\}}\theta_u^n\leq C[\capacity_{v_2}(u<v-s)]^{1+\alpha}.
		\end{equation}
		\end{itemize}
	Then, for every $\varepsilon>0$,
	\begin{equation*}
		\sup_{X }(v-u)\leq \varepsilon+\dfrac{2MC^{1/n}}{1-2^{-\alpha}}[\capacity_{v_2}(u<v-\varepsilon)]^{\alpha/n}.
	\end{equation*}
\end{lemma}
\begin{proof} {We adapt the same arguments as in~\cite{eyssidieux2009singular}.}
	By Lemma \ref{lem t^n cap}, for every $s>0$ and $0\leq t\leq 1-\lambda$, we have
		$$t^n \capacity_{v_2}\{u<v-s-t\}\leq \int_{\{u<v-s\}}\theta_u^n.$$
		Therefore, by \eqref{eq0lemsup<cap}, we obtain
		$$t^n \capacity_{v_2}\{u<v-s-t\}\leq C[\capacity_{v_2}(u<v-s)]^{1+\alpha},$$
		for every $s>0$ and $0\leq t\leq 1-\lambda$. Set $g(s)\coloneqq\left[\capacity_{v_2}(u<v-s)\right]^{1/n}$. We have
		\begin{equation*}
			t g(t+s)\leq C^{1/n}  g(s)^{1+\alpha},
		\end{equation*}
	for every $s>0$ and $0\leq t\leq 1-\lambda$.	
	If there exists $0\leq t_0\leq 1$ such that $g(\varepsilon)^{\alpha}=\frac{t_0}{2 C^{1/n}}$, then
	it follows from Lemma \ref{lem tg(t+s)} that $g\left(\varepsilon+\frac{t_0}{1-2^{-\alpha}}\right)=0$. Hence
	$$\sup_X (v-u)\leq \varepsilon+\dfrac{t_0}{1-2^{-\alpha}}\leq\varepsilon+\dfrac{2 C^{1/n}g(\varepsilon)^{\alpha}}{1-2^{-\alpha}},$$
    since $1\leq M$.
	Otherwise $2C^{1/n}g(\varepsilon)^{\alpha}\geq 1$, we infer that
		$$\sup_X(v-u)\leq M 
		\leq\varepsilon+\dfrac{2 M C^{1/n}}{1-2^{-\alpha}}g(\varepsilon)^{\alpha}.$$
This completes the proof.
\end{proof}
\begin{lemma}\label{lem extend GZ12 prop 5.3} Under the assumption of Lemma \ref{lem sup<cap}, assume further that $v_2\geq u$. We have
the following
	inequality 
$$\sup_{X } (v_1-u)\leq \left(2+\dfrac{2 M C^{1/n}}{\lambda^{1+\alpha}(1-2^{-\alpha})}\right)\|(v_1-u)_+\|_{L^1(\nu)}^{\frac{\alpha}{(n+1)\alpha+n}},$$
where $\nu=\mathbf{1}_{\{u<v_1\}}\theta_u^n$ and $(v_1-u)_+=\max(v_1-u,0)$.
\end{lemma}
\begin{proof}
	It follows from Lemma \ref{lem sup<cap} that for every $\varepsilon>0$, we have
	\begin{equation}\label{eq: lem35.1}
	\sup_X(v-u)\leq 2\varepsilon\lambda+\dfrac{2 M C^{1/n}}{1-2^{-\alpha}}[\capacity_{v_2}(u<v-2\varepsilon\lambda)]^{\alpha/n}.
    \end{equation}
	By Lemma \ref{lem t^n cap}, we obtain
		\begin{equation}\label{eq: lem35.2}
		    \begin{split}
		        (\varepsilon\lambda)^n\capacity_{v_2}(u<v-2\varepsilon\lambda)\leq \int_{\{u<v-\varepsilon\lambda\}}\theta_u^n\leq \int_{\{u<v_1-\varepsilon\}}\theta_u^n
		    \end{split}
		\end{equation} because $u\leq v_2$ so $\{u<v-\varepsilon\lambda\}\subset\{u<v_1-\varepsilon\}$. 
   Plugging \eqref{eq: lem35.2} into \eqref{eq: lem35.1}, since $v_2\geq u$ we obtain     
        \begin{align*}
		    \sup_X(v_1-u)&\leq 
2\varepsilon+\dfrac{2 M C^{1/n}}{\varepsilon^{\alpha}\lambda^{1+\alpha} (1-2^{-\alpha})}\left(\int_{\{u<v_1-\varepsilon\}}\theta_u^n\right)^{\alpha/n}
\\
&\leq 2\varepsilon+\dfrac{2 M C^{1/n}}{\varepsilon^{\alpha}\lambda^{1+\alpha} (1-2^{-\alpha}) }\left(\frac{1}{\varepsilon}\int_{\{u<v_1-\varepsilon\}}(v_1-u)_+d\nu\right)^{\alpha/n} \\
&\leq  2\varepsilon+\dfrac{2 M C^{1/n}}{\varepsilon^{\alpha(n+1)/n}\lambda^{1+\alpha}(1-2^{-\alpha}) }\|(v_1-u)_+\|_{L^1(\nu)}^{\alpha/n},
		\end{align*} 
        where the second inequality is due to Chebyshev’s inequality.
Choosing $\varepsilon:=\|(v_1-u)_+\|_{L^1(\nu)}^{\frac{\alpha}{(n+1)\alpha+n}}$ yields the desired inequality.
\end{proof}
% We now prove Theorem~\ref{main}. 
\begin{proof}[Proof of Theorem~\ref{main}] 
	In this proof, we denote by	$M_j(a)$, for $j\geq 1$, 
    constants in $[1, \infty)$ only depending on $n, X, \omega_X, \theta, B, \beta$ and an upper bound for $\int_X e^{\frac{2(V_\theta-u)}{a\varepsilon_0}} d\mu$. We assume that $a\mapsto M_j(a)$ is decreasing for every
	$j\geq 1$, and $\lim_{a\to 0^+}M_{j}(a)=\infty$.

\smallskip
Since the cohomology class $\{\theta\}$ is big, there exists a negative 
$\theta$-psh function $\Psi_0$ with analytic
singularities such that $\theta +dd^c\Psi_0\geq\varepsilon_0\omega_X$
for some $\varepsilon_0>0$. 
It follows from~\cite{boucksom2004divisorial} that we can choose $\Psi_0$ so that ${\rm Amp}(\theta)=X\setminus {\rm Sing}(\Psi_0)$. By subtracting a positive constant, we can assume $\Psi_0\leq V_\theta$.

Fix $0<c_1\leq \frac{\varepsilon_0}{4A}$ 
and $0<\delta_1<\min\{\delta_0, \sqrt{\frac{\varepsilon_0}{4K}}\}$, where $A$, $K$, and $\delta_0$ are defined as in Lemma~\ref{lem kiselman}. For $0<\delta<\delta_1$,
we define
\begin{equation}\label{eq: constantB0}
B_0=B_0(c, \delta)=\frac{Ac+K^2\delta}{\varepsilon_0},
\end{equation}
with $c=c(\delta)$ being chosen hereafter. We set
\begin{equation}\label{eq: u c delta}
u_{c, \delta}:=B_0\Psi_0+(1-B_0)\Phi_{c, \delta},
\end{equation}
where $\Phi_{c, \delta}$ is defined by \eqref{kisleg}. It follows from Lemma~\ref{lem kiselman} that $$\theta+\ddc u_{c,\delta} \geq \varepsilon_0B_0^2\omega_X,$$ hence $u_{c,\delta}$ is $\theta$-psh. 
% From now on, we set $\Omega:={\rm Amp}(\theta)\setminus\{\psi=-\infty\}$.
% (Use Lemma \ref{lem DL17}, Lemma \ref{lem extend GZ12 prop 5.3},
% Lemma \ref{Jen} and Lemma \ref{lem thmD Dem et al}.) 
% Assume $0<\delta<1$. We define $B_0$ and $u_{c,\delta}$ as in \eqref{eq: constantB0} and \eqref{eq: u c delta}:
% $$B_0=B_0(c, \delta)=\frac{Ac+2K\delta}{\varepsilon_0},$$
% and
% $$u_{c, \delta}=2B_0\Psi_0+(1-2B_0)\Phi_{c, \delta},$$
% where $c=c(\delta)>0$ is a small constant that will be chosen later.
% Note that we can choose $c$ such that $\frac{c}{\delta}\rightarrow\infty$ as $\delta\searrow 0$. 
% By Lemma \ref{kisleg}, we have $$\theta+dd^c u_{c, \delta}\geq \varepsilon_0B_0\omega_X.$$
Recall that $a_0$ is a positive constant such that $\omega_X+a_0\ddc\psi\geq 0$. Fixing $0<a\leq a_0$, we set $$v_{a, \delta}:=u_{c, \delta}+ a\varepsilon_0B_0^2\psi.$$
 % We now divide into several steps. 
\smallskip
 \textbf{Step 1: Bounding $\sup_X(v_{a,\delta}-u)$}. 
  We observe that $v_{a, \delta}$ is $\theta$-psh and $v_{a, \delta}\leq V_{\theta}+a \varepsilon_0B_0^2\psi.$ 
 Applying Lemma~\ref{lem DL17} with $\lambda=B_0$ and $v=v_{a, \delta}$, we obtain 
 \begin{equation}\label{eq1main}
 u\geq B_0v_{a, \delta}+(1-B_0)V_\theta -M_1(aB_0^2)\geq v_{a, \delta}-M_1(aB_0^2),
 \end{equation} using that $v_{a,\delta}\leq V_\theta$. 
 Since $v_{a, \delta}\leq V_{\theta}+aB_0^2 \varepsilon_0\psi,$ we also have $-\psi\leq \frac{V_{\theta}-u}{aB_0^2\varepsilon_0}$
 on $\{u<v_{a, \delta}\}$. We infer that for $s\geq 0$,
 \begin{align*}
\int_{\{u<v_{a, \delta}-s\}}\theta_{u}^n=\int_{\{u<v_{a, \delta}-s\}}e^{-\psi}d\mu 
		&\leq \int_{\{u<v_{a, \delta}-s\}}e^{\frac{V_{\theta}-u}{aB_0\varepsilon_0}}d\mu\\
&\leq \left(\int_{X}e^{\frac{2(V_{\theta}-u)}{aB_0^2\varepsilon_0}}d\mu\right)^{1/2} \left(\int_{\{u<v_{a, \delta}-s\}}d\mu\right)^{1/2}\\
&\leq M_2(aB_0^2)\left[\capacity_{V_{\theta}}(\{u<v_{a, \delta}-s\})\right]^{2},
 \end{align*}
 where the second inequality follows from H\"older's inequality and the last one follows from Theorem \ref{the exp cap}.
%  On $\{u<v_{a, \delta}\}$, we have
%  \begin{align*}
%      v_{a/2,\delta}+\frac{aB_0^2\varepsilon_0}{2} \psi=v_{a,\delta}>u,
%  \end{align*}
% so \begin{align*}
%     -\psi\leq \frac{2}{aB_0^2\varepsilon_0}(v_{a/2,\delta}-u)\leq  \frac{2M_1\big(\frac{aB_0^2}{2} \big)}{aB_0^2\varepsilon_0}
% \end{align*} using ~\eqref{eq1main}.
  Combining this with Lemma \ref{lem extend GZ12 prop 5.3} ($\lambda=B_0$, $v_1=v_{a,\delta}$, $v_2=V_\theta$), we infer that
\begin{equation}\label{eq2main}
\begin{split}
    \sup_X (v_{a, \delta}-u)&\leq \frac{M_3(aB_0^2)}{B_0^2}\|v_{a, \delta}-u\|^{\frac{1}{2n+1}}_{L^1(\nu)}\\
&\leq \frac{M_4(aB_0^2)}{B_0^2}\|(v_{a, \delta}-u)_+\|^{\frac{1}{2n+1}}_{L^1(\mu)},
\end{split}
\end{equation}
for every $s\geq 0$, where $\nu=\mathbf{1}_{\{u<v_{a, \delta}\}}\theta_u^n$.
Applying Lemma~\ref{lem DL17} with $\lambda=\frac{1}{2}$ and $v=\Psi_0+\varepsilon_0a\psi$, we obtain 
\begin{equation}\label{eq0.1main}
	u\geq \dfrac{\Psi_0+\varepsilon_0a\psi+V_{\theta}}{2}-M_5(a)\geq \Psi_0+\dfrac{\varepsilon_0a\psi}{2}-M_5(a).
\end{equation}
It follows that
\begin{align*}
v_{a, \delta}-u&=2B_0(\Psi_0+\dfrac{\varepsilon_0a\psi}{2}-u)+(1-B_0)(\Phi_{c, \delta}-u)\\
&\leq 2B_0 M_5(a)+(\Phi_{c, \delta}-u)_+.
\end{align*}
This combined with \eqref{eq2main} implies that

\begin{equation}\label{bound v a delta}
    \begin{split}
         \sup_X (v_{a, \delta}-u)&\leq \sup_X (v_{a, \delta}-u-2B_0 M_5(a))+2B_0 M_5(a)\\
    &\leq\frac{M_4(aB_0^2)}{B_0^2}\|(v_{a, \delta}-u-2B_0 M_5(a))_+\|^{\frac{1}{2n+1}}_{L^1(\mu)}+2B_0 M_5(a) \\
    &\leq\frac{M_4(aB_0^2)}{B_0^2}\|(\Phi_{c, \delta}-u)_+\|^{\frac{1}{2n+1}}_{L^1(\mu)}+2B_0 M_5(a) \\
      &\leq\frac{M_6(aB_0^2)}{B_0^2}\|(\Phi_{c, \delta}-u)_+\|^{\frac{\beta}{2n+1}}_{L^1(\omega_X^n)}+2B_0 M_5(a) \\
   &\leq\frac{M_6(aB_0^2)}{B_0^2}\|(\rho_{\delta}u-u)_+\|^{\frac{\beta}{2n+1}}_{L^1(\omega_X^n)}+2B_0 M_5(a)\\
    &\leq\frac{M_7(aB_0^2)}{B_0^2}\delta^{\frac{2\beta}{2n+1}}+2B_0 M_5(a),
    \end{split}
\end{equation}
where the last estimate holds due to Lemma \ref{Jen}.
Fix $U\Subset \textrm{Amp}(\theta)\setminus\{
\psi=-\infty\}$. 

\smallskip
\noindent\textbf{Step 2: Bounding $\sup_U(u_{c,\delta}-u)$.}
Recall that $u_{c,\delta}=v_{a,\delta}-aB_0^2\varepsilon_0\psi$. Set $m_U:=\sup_{\Bar{U}}(-\psi)$.
It follows from \eqref{bound v a delta} that
$$\sup_U(u_{c, \delta}-u)\leq aB_0^2\varepsilon_0m_U+\frac{M_7(aB_0^2)}{B_0}\delta^{\frac{2\beta}{2n+1}}+2B_0 M_5(a).$$
Denoting by $h_1$ the function $a\mapsto a\varepsilon_0m_U+2M_5(a)$ and choosing $$a\coloneqq\min\left\{h_1^{-1}\left(\frac{1}{\sqrt{B_0}}\right),\frac{a_0}{2}\right\},$$ since $B_0^2\leq B_0$ we have
$$\sup_U(u_{c, \delta}-u)\leq \sqrt{B_0}+\frac{M_7\left(h_1^{-1}\left(\frac{1}{\sqrt{B_0}}\right)B_0^2\right)}{B_0^2}\delta^{\frac{2\beta}{2n+1}}.$$
Denoting by $h_2$ the function $t\mapsto \frac{M_7\left(h_1^{-1}\left(\frac{1}{\sqrt{t}}\right)t^2\right)}{\sqrt{t^5}}$ and choosing $c=c(\delta)>0$ so that
$$h_2(B_0)=h\left( \frac{Ac+K\delta^2}{\varepsilon_0}\right)=\delta^{\frac{-2\beta}{2n+1}}.$$ 
Thus, we obtain
$$\sup_U(u_{c, \delta}-u)\leq 2\sqrt{B_0(\delta)} .$$
\noindent \textbf{Step 3: Conclusion.}
As follows from Lemma \ref{lem thmD Dem et al} below, we obtain
\begin{equation*}
    \rho_{\kappa(\delta)}u(z)-u(z)\leq 4\sqrt{B_0(\delta)}+2C_U B_0(\delta)+2K^2\delta,
\end{equation*} where $$\kappa(\delta)=\delta\exp\Big(\dfrac{-4A(C_UB_0+2\sqrt{B_0(\delta)}+2K^2\delta)}{\varepsilon_0B_0} \Big),\quad C_U=\sup_U(V_\theta-\Psi_0).$$
% {\color{red} Is $\kappa$ constant, not depending on $\delta$???}
Replacing $\delta$ with $\kappa (\delta)$ we obtain for $\delta\leq \kappa(\delta_0)$,
\[\rho_\delta u(z)-u(z)\leq 4\sqrt{B_0(\kappa^{-1} (\delta))}+2C_U B_0(\kappa^{-1} (\delta)). \]
% {\color{red} We check that there exists $C>1$ such that \[ \limsup_{t\to 0^+}\Bigg(\frac{g(At)}{Ag(t)}\Bigg)<\frac{1}{2n}.\]}
The conclusion thus follows.
\end{proof}

\begin{lemma}\label{lem thmD Dem et al}
	Let $c: (0, \delta_1)\rightarrow (0, c_1)$ and $B_1:\R^+\rightarrow\R^+$ be functions. Assume $\lim_{\delta\to 0^+}c(\delta)/\delta=+\infty$, so that $K^2\delta<Ac(\delta)$ for every $0<\delta<\delta_2$, where $\delta_2\in (0, \delta_1)$ is sufficiently small. 
    Let $u\in\PSH(X,\theta)$ be such that
	\begin{equation}\label{eq0 lem thmD Dem et al}
	\sup_{U}(u_{c(\delta), \delta}-u)\leq  B_1(\delta),
	\end{equation}
	for every $\delta\in(0,\delta_2)$, where $u_{c(\delta),\delta}$ is defined by \eqref{eq: u c delta} and 
    $U\Subset \textrm{Amp}(\theta)\backslash\{\psi=-\infty\}$. Then we have \begin{equation}\label{eq: delta_diff}
	\sup_U(\rho_{\kappa(\delta)}u-u)\leq 2(B_1+C_UB_0+2K^2\delta),
	\end{equation}
where
	\begin{equation*}
	    C_U=\sup_U(V_{\theta}-\Psi_0) \;\text{and}\; \kappa(\delta)=\delta\exp\Big(\dfrac{-4A(C_UB_0+B_1+2K^2\delta)}{\varepsilon_0B_0} \Big).
	\end{equation*}
\end{lemma}
\begin{proof}The argument follows from \cite[pages 642-643]{demailly2014holder}, we include the proof
for the readers’ convenience.
	By the assumption, we have
\begin{align*}
B_1\geq u_{c, \delta}(z)-u(z)&=B_0(\Psi_0(z)-u(z))\\
&+
(1-B_0)\left(\rho_{t_0}u(z)+K(t_0^2-\delta^2)+K(t_0-\delta)-c\log\left(\frac{t_0}{\delta}\right)-u(z)\right),	
\end{align*}
for every $z\in U$, where $t_0=t_0(z)\in (0, \delta)$ is a minimum point of 
$\rho_{t}u(z)+Kt^2+Kt-c\log\left(\frac{t}{\delta}\right)$. Since $V_{\theta}\geq u$
and $\rho_{t_0}u+Kt_0^2\geq u$ (see Lemma \ref{lem kiselman}), it follows that
$$B_0(\Psi_0(z)-V_{\theta}(z))
-c(1-B_0)\log\left(\frac{t_0}{\delta}\right)\leq  B_1+K\delta^2+K\delta\leq B_1+2K^2\delta,$$
for every $z\in U$. Hence, for every $z\in U$,
\begin{align*}
t_0(z)&\geq \delta\exp\bigg(\dfrac{B_0(\Psi_0(z)-V_{\theta}(z))-B_1-2K^2\delta}{c(1-B_0)} \bigg)\\
&\geq \delta\exp\Big(\dfrac{-4C_UB_0-4B_1-8K^2\delta}{A^{-1}\varepsilon_0B_0} \Big)
= \kappa(\delta),
\end{align*}
 using the facts that $B_0\varepsilon_0=Ac+K^2\delta\leq 2Ac$ and $1-B_0\geq \frac{1}{2}$. Since $t\mapsto \rho_tu+Kt^2$ is increasing, we have
\begin{align*}
    \rho_{\kappa(\delta)}u(z)-u(z)
	&\leq \rho_{t_0}u(z)+Kt_0^2-u(z)\\
		&\leq \left(\rho_{t_0}u(z)+K(t_0^2-\delta^2)+K(t_0-\delta)-c\log\left(\frac{t_0}{\delta}\right)\right)-u(z)+2K^2\delta\\
	&\leq \Phi_{c,\delta}(z)-u(z)+2K^2\delta\\
	% &=\dfrac{1}{1-B_0}(u_{c, \delta}(z)-B_0\Psi_0(z))-u(z)\\
	&=\dfrac{1}{1-B_0}(u_{c, \delta}(z)-u(z))+\dfrac{B_0}{1-B_0}(u(z)-\Psi_0(z))+2K^2\delta\\
	&\leq \dfrac{B_1(\delta)}{1-B_0}+\dfrac{B_0}{1-B_0}(V_{\theta}(z)-\Psi_0(z))+2K^2\delta,
\end{align*}	
for every $z\in U$, where the last inequality holds due to \eqref{eq0 lem thmD Dem et al}. Therefore,
\begin{center}
	$\rho_{\kappa(\delta)}u(z)-u(z)\leq \dfrac{B_1}{1-B_0}+\dfrac{C_UB_0}{1-B_0}+2K^2\delta\leq 2(B_1+B_0C_U)+2K^2\delta,$
\end{center}
for every $z\in U$.
The proof is thus complete.
\end{proof}

We now deal with the special case when $\theta =\omega_X$ is a K\"ahler form, proving Corollary~\ref{coro: main}. 
Recall that $\mathcal{N}=\mathcal{N}(B,\beta,a_0,C_0,h)$ is the set of probability measures $\nu$ on $X$ such that $\nu=e^{-\psi}\mu$ with $\mu\in\mathcal{M}(B,\beta)$, $\psi\in\PSH(X, a_0\omega_X)$ and  
$$\int_Xh(-\psi)e^{-\psi}{\rm d}\mu\leq C_0,$$ 
where $a_0>0$, $ C_0>0$ are fixed constants and $h: (0, \infty)\rightarrow (0, \infty)$ is an increasing concave function with $h(\infty)=\infty$. We define the set $$\mathcal{F}=\{u\in\mathcal{E}(X,\omega_X),\,\sup_X u=0\colon (\omega_X+\ddc u)^n\in\mathcal{N}\}.$$
\begin{lemma}\label{lem: compact}
   The closure of $\mathcal{F}$ is contained in $\mathcal{E}(X,\omega_X)$. 
\end{lemma}
\begin{proof} Let $(u_j)_{j\in\mathbb N}\subset \mathcal{F}$ be a sequence of $\omega_X$-psh functions such that $(\omega_X+\ddc u_j)^n\in\mathcal{N}$ and $\sup_X u_j=0$. By Hartogs' lemma (see, e.g., \cite[Proposition 8.4]{guedj2017degenerate}), there exists a subsequence, still denoted $(u_j)$, such that
$u_j\xrightarrow{L^1} u$ for some $u\in\PSH(X, \omega_X)$. 
    By Proposition~\ref{prop: unif-capa}, the family of non-pluripolar measures $(\omega_X+\ddc u_j)^n$ is uniformly absolutely continuous with respect to capacity. Set $v_j=\max(u_j,u)$. We see that $v_j$ converges to $u$ in capacity thanks to Hartogs' lemma and the quasi-continuity of $\omega_X$-psh functions.
    As follows from~\cite[Propositions 2.10]{dinew12-cv-capacity}, the sequence of measures $(\omega_X+\ddc v_j)^n$ is uniformly absolutely continuous with respect to capacity. Applying~\cite[Proposition 2.11]{dinew12-cv-capacity}, we conclude that $u\in\mathcal{E}(X,\omega_X)$.
    Therefore, the closure of $\mathcal{F}$ in the $L^1$ topology is contained in $\mathcal{E}(X,\omega_X)$.
\end{proof}
\begin{proof}[Proof of Corollary~\ref{coro: main}] 
    By Theorem~\ref{main}, it suffices to show that there exists a positive constant $C=C(a_0,X,\omega_X,\mu)$ such that
    \begin{equation}\label{eq: skoda}
        \int_X e^{-2u/a_0}d\mu\leq C.
    \end{equation}
    Thanks to Lemma~\ref{lem: compact}, the closure of the set $\mathcal{F}= \{u\in\mathcal{E}(X,\omega_X),\,\sup_X u=0\colon (\omega_X+\ddc u)^n\in\mathcal{N}\}$ is a compact family of functions in $\mathcal{E}(X,\omega_X)$. Since functions in $\mathcal{E}(X,\omega_X)$ have zero Lelong number at every point on $X$ by \cite{guedj2007weighted}, the uniform Skoda integrability
theorem (see, e.g.,~\cite{guedj2017degenerate}) ensures that there exists a constant $C$ depending on $a_0$ and $\overline{\mathcal{F}}$ such that the inequality \eqref{eq: skoda} holds for all $u\in\mathcal{F}$. 
\end{proof}
\begin{remark}
     % Lemma~\ref{lem: compact} still holds in the more general setting when $\theta$ is merely semi-positive and big (i.e., $\int_X\theta^n>0$), so does Corollary~\ref{coro: main} but the modulus of continuity is only carried on $\textrm{Amp}(\theta)\backslash\{\psi=-\infty\}$. 
     Lemma~\ref{lem: compact} still holds in the more general setting where $\theta$ is merely semi-positive and big (i.e., $\int_X \theta^n > 0$). Corollary \ref{coro: main} also remains valid, but in this case only on compact subsets of $\textrm{Amp}(\theta) \setminus \{\psi = -\infty\}$.
     % Hence, the choice of the function $F_U$ in Theorem~\ref{main} does not depend on the integral $\int_X e^{-2u/a_0}d\mu$.
\end{remark}
\bibliographystyle{alpha}
	\bibliography{bibfile}	

\end{document}
